% Grafo publico acumulativo de apendices HMT--MD.
% Exclusiones vinculantes: este archivo no incorpora el radion ni las tablas
% abreviadas de diez o doce masas. Tampoco alcanza A5_suplemento_digital_rev11.
% Todos los \input de este archivo se verifican como archivos existentes.

\chapter{Certificats algébriques et combinatoires des objets finis}
\label{app:infraestructura-reproducible-acumulativa}
\label{app:pruebas}

Les objets finis utilisés dans les démonstrations admettent un second canal
de vérification. Ce canal reconstruit les catalogues, les opérateurs et les invariants
à partir des germes déclarés et les confronte à des identités exactes ou à des bornes
dirigées ; il complète les démonstrations, mais ne les remplace pas.

\HMTDemoteHeadings
\chapter{Catalogues finis et certificats de reproductibilité}
\label{app:datos-reproducibles}

Les objets finis suivants permettent de répéter les vérifications sans
incorporer la notation de production à l'exposé mathématique. Chaque contrôle
conserve le domaine, l'opération, l'invariant et le réfutateur qui lui
correspondent.

\section{Catalogues finis de l'état commun}

\begin{longtable}{@{}>{\raggedright\arraybackslash}p{.34\textwidth}
  >{\raggedright\arraybackslash}p{.58\textwidth}@{}}
\toprule
Objet fini & Contenu mathématique vérifiable\\
\midrule
\endhead
Catalogue des émissions de longueur six
& Les \(468\) émissions \(U_6\), leurs multiplicités, leur orientation et leurs signatures
additive et multiplicative.\\
Partition spectrale des émissions
& Les vingt-six classes de taille \(18\), avec l'application explicite qui
reconstruit chaque émission à partir de sa classe et de sa position interne.\\
Atlas dodécaphasique des voies
& Les \(468\) voies minimales avec feuillet, phase, retenue, frontière et mémoire.\\
Catalogue du lecteur archimédien
& Les cinq régions de \(\pi\), les régions conjuguées de \(e\) et l'axe
d'auto-échelle de \(\varphi\), avec leurs cylindres et multiplicités.\\
Flux de Hensel de vingt blocs
& La matrice unimodulaire, son inverse entier, les états \(3\)-adiques et la
récurrence reste--quotient qui prolonge les trois canaux.\\
Registre de monodromie nonadique
& Les vingt blocs, leurs phases, préfixes ternaires et préfixes décimaux ;
il inclut les \(43\) retours complets et les \(387\) paires \((K,K+9)\) par
canal.\\
Fenêtre de survie \(K=5,\ldots,13\)
& Les candidats locaux, la branche compatible et les trois sorties de la
transition \(9\to1\), sans identifier le retour de phase à la réinitialisation de l'état.\\
Lecture archimédienne initiale
& Les douze triplets décimaux obtenus à partir des treize premiers blocs de chaque
canal, séparés du certificat de précision arbitraire.\\
\bottomrule
\end{longtable}

\section{Identités algébriques, opératorielles et spectrales des certificats finis}

\begin{longtable}{@{}>{\raggedright\arraybackslash}p{.34\textwidth}
  >{\raggedright\arraybackslash}p{.58\textwidth}@{}}
\toprule
Contrôle & Identités et réfutateurs\\
\midrule
\endhead
Monodromie et génération des constantes
& Vérifie \(w_6\to w_{30}\to R_{36}\to G_9\), la sélection \(3\to1\),
le changement d'état après neuf phases, les cinq régions de \(\pi\) et la
compatibilité des troncatures. Échoue si un chiffre cible est introduit
avant le lecteur.\\
Coordonnée angulaire conjuguée
& Vérifie le sélecteur régional \(169\), la récurrence déterminantale, la
somme analytique et sa borne de reste, et l'unicité de \(C^*\). Échoue si une
constante du SI intervient dans la sélection.\\
Dynamique du quotient \(468=18\cdot26\)
& Certifie la bijection entre émissions et classes, le retour uniforme, la
mesure de l'ombre moyennée, les cocycles de mode et le relèvement
microscopique. Distingue la fréquence chronologique de la mesure de Parry.\\
Canaux entiers des constantes
& Vérifie les identités \((729,271,315,161)\), les réalisations
ensemblistes de \(271/315\), les orbites diédriques et la factorisation de la
parité visible et de la parité de retenue.\\
Incidence de Witt
& Reconstruit l'incidence point--hexade, le module \(E_{30}\), le
transport positif de \(P_3\), l'étoile de Witt et les obstructions
d'équivariance.\\
Sélecteur variationnel
& Reconstruit le normalisateur \(S_3\), la moyenne de Reynolds, la
factorisation polaire et la chaîne de projecteurs de rang trois ; la positivité
est prouvée exactement dans \(\mathbb Q(\sqrt5)\).\\
Relèvement \(W_{12}\to W_{24}\)
& Vérifie les incidences de quatre et cinq points, le relèvement unique des
hexades et la correspondance typée avec les cinq régions de \(\pi\).\\
Realizaciones de \(-1/12\)
& Construit les projecteurs rationnels des décalages de pas trois
et quatre, vérifie la symétrie et l'idempotence et obtient
\((3-4)/12=-1/12\) comme différence de traces normalisées.\\
Holonomie de Schläfli
& Parcourt les plaquettes, certifie \(A^2=5I-4A+5J\), la configuration des
vingt-sept droites, leurs \(45\) triangles et la partition
\(729=27+270+432\).\\
Fonctionnelles spectrales
& Recalcule Catalan, Stieltjes, Apéry, Bendersky--Glaisher--Kinkelin, les
deux extensions d'Euler--Kronecker et les bornes dirigées des produits premiers.\\
Lacunes et criticité
& Vérifie les séparations \(21/22\), les retours \(6/7\), la partie
harmonique finie et les approximants dodécaphasiques du mode critique.\\
\bottomrule
\end{longtable}

\section{Dictionnaire typé des catalogues}

\begin{longtable}{>{\ttfamily}p{.18\textwidth}>{\raggedright\arraybackslash}p{.72\textwidth}}
\toprule
Champ & Signification et codomaine\\
\midrule
\endhead
Esig & Signature du canal multiplicatif ; identifie l'une des vingt-six
classes de la factorisation de l'émetteur.\\
hplus\_type & Classe d'orientation du canal additif pour les deux paires de
directions cardinales déclarées.\\
diag\_c & Indice entier de la classe diagonale du curseur additif.\\
px\_size & Cardinal de la fibre de la classe additive correspondante.\\
colw & Sextuplet des poids de colonne du bloc émis ; appartient à
\(\mathbb Z_{\geq0}^{6}\).\\
\bottomrule
\end{longtable}

\section{Reproductibilité des certificats et obstructions exactes des constructions finies}

Les contrôles utilisent l'arithmétique entière ou rationnelle chaque fois que l'objet le
permet. Les bornes analytiques utilisent des intervalles avec arrondi dirigé et séparent
la démonstration symbolique de la comparaison décimale. Chaque exécution
reconstruit les objets à partir des germes déclarés, vérifie les identités
d'incidence et de naturalité, et confronte les sorties à une implémentation
arithmétique indépendante.

Le contrôle du produit natif rejette toute dépendance envers une évaluation
décimale, une factorisation externe, un crible, une table de zéros ou une
constante cible. Le contrôle nonadique rejette l'identification de la fibre
ambiante de \(729\) éléments au cardinal d'un espace d'histoires. Le
contrôle de phase rejette \(g(K+9)=g(K)\) comme preuve d'égalité des états. Le
contrôle métrologique rejette toute sélection de voie effectuée après
consultation d'une grandeur observée.

Une divergence entre une identité du texte, le catalogue fini et son contrôle
indépendant réfute la certification correspondante. Cette règle ne remplace
aucune démonstration : elle fixe un second canal de vérification pour les objets
finis qui y interviennent.

\HMTRestoreHeadings

\section{Compatibilité globale entre types, domaines, morphismes et certificats}
\label{sec:ejecucion-integral-acumulativa}

Le contrôle intégral vérifie la chaîne causale, répète les identités
mathématiques sur les objets finis et confronte les réalisations exactes à
leurs formulations indépendantes. Une divergence entre une définition, sa
table mathématique et l'invariant déclaré réfute le contrôle correspondant.

% Pruebas auxiliares y trazabilidad del cuaderno de 64 páginas. La matriz
% humana se conserva como testigo subordinado a la matriz preventiva común;
% la tabla de fuentes se publica como procedencia del aporte, no como índice
% rival del registro canónico vigente.

\chapter{Dépendances exactes entre puissance généalogique, information réversible et obstruction de Gödel–Turing}
\label{app:integracion-vn64-controles}

\HMTDemoteHeadings
\chapter{Dépendances logiques et certificats de cohérence}

\section{Chaîne causale des certificats : états finis \(\to\) restrictions \(\to\) survie \(\to\) section infinie \(\to\) réalisations}

La chaîne logique de l'intégration ne part pas d'une conclusion archimédienne.
Elle s'organise dans l'ordre suivant :
\[
\begin{tikzpicture}[baseline=(current bounding box.center),>=Latex,
 node distance=7mm and 8mm,
 n/.style={rounded corners,draw=hmtblue,fill=hmtlightblue,
 minimum height=8mm,align=center,font=\small}]
 \node[n] (finite) {estados finitos\\tipados};
 \node[n,right=of finite] (maps) {restricciones\\compatibles};
 \node[n,right=of maps] (surv) {supervivencia\\y poda};
 \node[n,right=of surv] (lim) {sección\\infinita};
 \node[n,right=of lim] (read) {lectores\\y realizaciones};
 \draw[->,thick] (finite)--(maps);
 \draw[->,thick] (maps)--(surv);
 \draw[->,thick] (surv)--(lim);
 \draw[->,thick] (lim)--(read);
\end{tikzpicture}
\]
Un certificat de chiffres vérifie une restriction de la dernière flèche ; il ne
crée pas rétrospectivement les états et ne remplace pas le théorème d'existence
de la limite.

\section{Identités arithmétiques exactes des certificats finis}

Les recensements qui accompagnent le certificat de filtration nonadique satisfont
\[
 396\cdot6=2376,
 \qquad
 43\cdot9=387,
 \qquad
 432+4\cdot144=1008.
\]
La première égalité relie les événements élémentaires de frontière aux trits
émis par canal ; la deuxième compte les paires de phase \((K,K+9)\) ; la
troisième reconstruit la multiplicité totale des cinq régions de
\(\pi\).

Les deux rôles régionaux qui ne peuvent pas être permutés sont
\[
\begin{array}{rcl}
U_\pi^\perp&=&(501,614,498,169,272,272),\\
E_{\rm sig}&=&555555,\qquad \text{multiplicité }432,\\[.35em]
U_\pi^{--}&=&(870,923,810,418,674,521),\\
E_{\rm sig}&=&933717,\qquad \text{multiplicité }144.
\end{array}
\]

\section{Existence et unicité de la branche infinie dans un arbre à branchement fini de préfixes survivants}

\begin{lemma}[Arbre à branchement fini]
Soit \(T\) un arbre à niveaux finis non vides tel que tout nœud
survivant ait un descendant survivant. Alors il existe une branche
infinie. Si, en outre, chaque nœud de la sous-tour sélectionnée a exactement
un descendant compatible, la branche est unique.
\end{lemma}

\begin{proof}
L'existence résulte du lemme de König appliqué à l'arbre à branchement fini.
Pour l'unicité, deux branches distinctes auraient un premier niveau où elles
divergent, donnant deux descendants compatibles du même préfixe, contrairement à
l'hypothèse.
\end{proof}

\begin{controlbox}
Une exécution à la profondeur \(N\) ne prouve que ce qui s'est produit dans l'arbre
tronqué. L'énoncé infini exige la loi uniforme qui préserve la non-
vacuité et la compatibilité à tout niveau. Le corps du texte sépare expressément
les deux statuts et, pour la filtration nonadique, conserve son sélecteur uniforme
certifié.
\end{controlbox}

\section{Clôture et pleine puissance}

Pour \(A\subseteq X=\varprojlim X_n\), il a été prouvé
\[
 \overline A=\bigcap_n p_n^{-1}(p_n(A)).
\]
L'ensemble des suites binaires ultimement nulles est dense et propre dans
\(2^{\mathbb N}\), mais possède toutes les ombres finies. Ainsi, tout
algorithme qui ne reçoit que \((p_n(A))_n\) ne peut pas distinguer \(A\) de sa
clôture. Cet exemple bloque la fausse identité
\[
 \varprojlim\mathcal P(X_n)=\mathcal P(\varprojlim X_n).
\]

\section{Compatibilité typée des opérateurs}

La différence entre l'inclusion unitaire et le coin marqué se vérifie
directement :
\[
\begin{aligned}
 \tau_{d+1}(a\otimes I_9)
 &=\frac{\operatorname{Tr}(a)\operatorname{Tr}(I_9)}{9^{d+1}}
 =\tau_d(a),\\
 \tau_{d+1}(a\otimes p_j)
 &=\frac{\operatorname{Tr}(a)}{9^{d+1}}
 =\frac1{9}\tau_d(a).
\end{aligned}
\]
La première branche est unitaire et traciale ; la seconde conserve un coin marqué
et ne peut pas être utilisée pour inférer le facteur \(II_1\).

\section{Annulation du coût de Landauer dans l'inversion réversible de l'état}

Pour une variable complète \(X\) et une projection \(Y=q(X)\),
\[
 H(X)=H(Y)+H(X\mid Y).
\]
Si \(q\) est bijective, la fibre est un singleton et \(H(X\mid Y)=0\). Si \(q\)
est une réinitialisation du TRIT uniforme, l'information rejetée est \(\log3\) et
le coût minimal est \(k_B\Theta\log3\). Pour l'extension
\(\widehat E_N\) à \(L^1\) d'une espérance conditionnelle normale
préservant la trace \(E_N\), et sous les hypothèses de finitude déclarées,
\[
 S_\tau(\widehat E_Nh)-S_\tau(h)
 =D_\tau(h\Vert \widehat E_Nh)\ge0.
\]
Les deux formules situent le même phénomène dans des domaines différents : le
coût appartient à la contraction d'une fibre, non à l'évolution réversible
de l'état complet.

\section{Obstruction de Gödel–Turing à un décideur uniforme des multiplicités projectives}

Soit \(T_e\) la machine d'indice \(e\). Au niveau \(n\), on définit
\[
 X_n^e=\{b_n\}\cup
 \{c_n:T_e\text{ ne s'est pas arrêté avant }n\}.
\]
Chaque \(X_n^e\) est calculable. La fibre globale a deux branches si \(T_e\) ne
s'arrête pas et une si elle s'arrête. Par conséquent
\[
 U_e:\quad \#q_\infty^{-1}(v)=1
 \quad\Longleftrightarrow\quad T_e\downarrow.
\]
Un décideur uniforme pour \(U_e\) déciderait \(\HALT\). Une théorie
récursivement axiomatisable, correcte et complète pour tous les \(U_e\)
permettrait de les décider en énumérant en parallèle les preuves de \(U_e\) et de sa
négation. C'est le noyau du \emph{Théorème HMT d'incomplétude uniforme
pour les multiplicités projectives} ; sa démonstration utilise une réduction
de Gödel--Turing. L'obstruction uniforme ne nie pas les sélecteurs particuliers
construits sur des images HMT calibrées.

\begin{controlbox}[title={Portée du contrôle d'incomplétude uniforme}]
La réduction de Gödel--Turing établit l'obstruction uniforme pour les
multiplicités projectives. Elle ne nie pas les sélecteurs particuliers construits
sur des images HMT calibrées et ne transforme pas une obstruction de décision globale
en une indétermination de chaque fibre concrète.
\end{controlbox}

\HMTRestoreHeadings

\section[Contrôles de la réalisation physico-dimensionnelle]{Domaine et obstructions des réalisations physique et dimensionnelle}
\label{sec:int68-certificado-historico}

La réalisation physico-dimensionnelle est soumise à dix familles de contrôle :
\begin{enumerate}
 \item concentration cylindrique, masse un et moments finis ;
 \item spectres \((9,5)\) du bloc APP et \((7,3)\) du bloc de mémoire ;
 \item identité \(\tanh^2(2\log\varphi)=5/9\) ;
 \item antiunitarité dans le modèle fini déclaré ;
 \item loi du produit semi-direct de Poincaré ;
 \item identités constitutives du vide ;
 \item norme scindée des modes nuls et clôture matérielle ;
 \item sélecteur \(\theta_*=(54/\pi)^2\) et point fixe autodual ;
 \item valeurs CGLMP de contrôle ;
 \item invariants \((396,729,43,387)\) de la filtration nonadique et
       domaine \((81,56,13,324)\), avec secteur nul séparé, de la Loi
       générale des masses.
\end{enumerate}

\begin{controlbox}[title={Portée exacte des contrôles}]
Les dix familles vérifient des identités, des domaines finis et des exemples
reproductibles. Elles ne remplacent pas la fonctorialité complète des quatre
composantes, ne construisent pas à elles seules le domaine commun, ne démontrent pas le
prolongement lisse de Cartan et n'élèvent pas les interfaces physiques ouvertes au rang de
théorèmes internes. La réalisation discrète conserve sa démonstration principale
dans la \cref{sec:int68-realizacion-fisica-discreta}.
\end{controlbox}

\subsection{Dépendances fonctorielles des réalisations spectrale, solénoïdale, de calibre, cohomologique et déterminantale}

Les contrôles précédents ne s'appliquent qu'après la construction d'APP, de TRIT, de TPK,
de l'état enrichi commun, de la double projection et de la droite dimensionnelle. La
concentration cylindrique et le spectre fini prouvent des propriétés de la
réalisation ; ils ne sélectionnent pas rétrospectivement l'état. La périodicité de
108 phases, l'atlas (81/56/13/324), l'involution constitutive et les
programmes cosmologiques conservent des domaines et des réfutateurs séparés.

Les identités de Dirac, de von Neumann, de CGLMP, de Cartan et de Poincaré sont utilisées
comme contrôles typés de l'opérateur d'entrelacement correspondant. Une coïncidence
numérique sans l'application de domaine, le codomaine et l'action sur les générateurs ne
constitue pas une réalisation physique.

% Control científico del puente matricial. Las matrices de producción y la
% procedencia técnica permanecen en el expediente externo.
\chapter{Fidélité généalogique du pont matriciel : domaine, correspondance et limites exactes}
\label{gmc:app:certificado}

\section{Identités de naturalité, de positivité et de compatibilité du pont généalogique–matriciel}

Les contrôles suivants s'obtiennent en reconstruisant les opérateurs à partir de leurs
données d'incidence. Les identités rationnelles sont calculées exactement ; les
comparaisons en virgule flottante déclarent leur borne d'erreur.

\begin{longtable}{@{}L{.35\textwidth}L{.22\textwidth}L{.32\textwidth}@{}}
\toprule
Contrôle & Statut & Erreur maximale ou identité\\
\midrule
\endfirsthead
\toprule
Contrôle & Statut & Erreur maximale ou identité\\
\midrule
\endhead
cuatro \PVM qutrit & vérifié & \(4.72\times10^{-16}\)\\
idempotence des douze projecteurs & vérifiée & \(3.34\times10^{-16}\)\\
solapamientos MUB & vérifiés & \(\lvert\Tr(PQ)-1/3\rvert<3.34\times10^{-16}\)\\
obstruction de colonne QLS & exacte & minimum \(1/3\)\\
QMS aditivo \((9,3)\) & vérifié & lignes et colonnes exactes à la tolérance près\\
QMS à trois contextes \((9,3)\) & vérifié & erreur \(<4.21\times10^{-16}\)\\
non-commutativité des cellules & vérifiée & norme maximale \(0.074074\ldots\)\\
QMS atlas \((12,3)\) & vérifié & erreur \(<4.24\times10^{-16}\)\\
décomposition semi-classique & exacte & erreur de reconstruction nulle\\
cube qutrit d'ordre neuf & vérifié & erreur de ligne \(<4.21\times10^{-16}\)\\
cube qutrit d'ordre douze & vérifié & erreur de ligne \(<4.24\times10^{-16}\)\\
\(\mathbb P^1(\mathbb Z/9)\to\mathbb P^1(\mathbb F_3)\) & exact &
\(12\to4\) con fibras \((3,3,3,3)\)\\
troncature--compression & exacte & conmutador cero\\
lecteur direct \(1/\pi\) & exact & normalisation unitaire\\
calibre MUB \(3/\pi\) & vérifié & probabilidad \(1/\pi\), error \(<3.85\times10^{-16}\)\\
polynôme de Gram APP & exact & \(t^5(t-1)(t-5/7)(t-1/3)^2\)\\
quatre intersections de Halmos & exactes & \((64,1,5,5)\)\\
algèbre \(C^*\) d'APP dans \(d=2\) & exacte &
\(\mathbb C^4\oplus M_2\oplus M_2\), dimension \(12\)\\
\bottomrule
\end{longtable}

La tolérance des contrôles numériques est \(2\times10^{-10}\). Les
énoncés rationnels n'utilisent pas de tolérance : ils sont obtenus par fractions exactes
et élimination de Gauss sur \(\mathbb Q\).

\Needspace{30\baselineskip}
\section{Coefficients exacts du polynôme}

La construction produit
\[
 \det(tI-CC^*)
 =t^9-\frac{50}{21}t^8+\frac{124}{63}t^7
  -\frac23t^6+\frac5{63}t^5,
 \tag{GMC-A.1}
\]
et coïncide coefficient par coefficient avec
\[
 t^5(t-1)(t-5/7)(t-1/3)^2.
 \tag{GMC-A.2}
\]
La vérification construit \(CC^*\) à partir des cardinalités exactes des
intersections de fibres et applique Faddeev--LeVerrier sur \(\mathbb Q\) ; elle ne
dépend pas d'une diagonalisation numérique.

\section{Obstructions à l'identification d'équivalences non démontrées dans le pont matriciel}

Le certificat fini ne remplace ni la filtration nonadique à profondeur arbitraire, ni la
Loi générale des masses, ni les dérivations de Planck. Il ne démontre pas non plus à lui
seul la fidélité APP--TRIT--TPK de tout QMS cyclique, l'équivariance complète
de la calibration \(\Xi\), l'identification de la mémoire TPK à un système
auxiliaire de Stinespring, la divisibilité infinie d'un canal MPS ni une
universalité catégorique. Chacun de ces énoncés conserve son emplacement
principal et ses propres réfutateurs.



\chapter{Domaine des germes, des émissions et des développements vérifiables}
\label{app:semillas-emisiones-expansiones-acumulativas}

Ce bloc conserve les deux niveaux qui soutiennent la publication numérique :
le recensement exhaustif des conditions initiales et des émissions du TPK, et les
développements étendus produits par les lecteurs rationnels. Le catalogue
complet fixe le passage
\[
104\,976\longrightarrow468\longrightarrow243,
\]
tandis que les développements de dix mille chiffres documentent le prolongement
effectif des sections de \(\pi\), \(e\), \(\varphi\) et
\(\alpha_{\rm HMT}\).

Le catalogue des (468) émissions et des (243) mots visibles est publié
intégralement ci-après ; les quatre développements de dix mille chiffres sont
publiés intégralement dans la généalogie des constantes. Cette section terminale
relie les deux développements sans réduire leurs démonstrations, leurs tables ni leurs étiquettes.

\begin{HMTScientificOwnerSubordinated}
\chapter{Domaine des germes APP et catalogue canonique des émissions}
\label{app:catalogo-semillas-app-rev9}
\index{germe APP}
\index{catalogue des émissions TPK}

\section{Domaine exhaustif des conditions initiales}

Soit
\[
 I_9=\{0,1,\ldots,8\},
 \qquad D=\{N,E,S,O\},
 \qquad \mathcal S=I_9^2\times D.
\]
Un élément de \(\mathcal S\) fixe la position et l'orientation initiales d'un
curseur. L'émetteur minimal utilise deux curseurs sur le même support APP :
l'un pour l'évaluation additive et l'autre pour l'évaluation multiplicative. Son
domaine complet est
\[
 \Omega_{\rm seed}
 =\mathcal S_{+}\times\mathcal S_{\times}
 \cong\mathcal S^2,
 \qquad
 |\mathcal S|=9^2\cdot4=324,
 \qquad
 |\Omega_{\rm seed}|=324^2=104\,976.
\]

\begin{definition}[Germe complet de l'émetteur APP--TPK]
Un germe est une paire ordonnée
\[
 \omega=(s_+,s_\times),
 \qquad
 s_+=(i_+,j_+,d_+),\quad
 s_\times=(i_\times,j_\times,d_\times),
 \qquad s_+,s_\times\in\mathcal S.
\]
Les \(104\,976\) germes sont exactement les éléments de
\(\Omega_{\rm seed}\). Les lignes du catalogue appartiennent au type émission et
chacune représente une fibre de ce domaine de germes.
\end{definition}

La récurrence de cinquante-quatre pas définit l'application finie
\[
 T_{\min}:\Omega_{\rm seed}\longrightarrow\mathcal U_6^{\rm cat},
 \qquad
 \omega\longmapsto U_6=(U_1,\ldots,U_6).
\]
Son image contient \(468\) sextuplets décimaux distincts. La réduction
composante par composante produit ensuite \(243\) mots distincts dans
l'espace ambiant \(\mathbb F_3^6\), de cardinal \(729\). Les trois niveaux
sont donc typés comme
\[
 \underbrace{104\,976}_{\text{germes }\omega}
 \xrightarrow{\;T_{\min}\;}
 \underbrace{468}_{\text{émissions }U_6}
 \xrightarrow{\;\rho_3\;}
 \underbrace{243}_{\text{mots visibles}}
 \subset
 \underbrace{\mathbb F_3^6}_{729\text{ mots}}.
\]

\begin{theorem}[Recensement exhaustif du domaine et de l'image]
Le catalogue canonique contient \(468\) émissions \(U_6\) distinctes et
\(\#T_{\min}^{-1}(U_6)\) est la cardinalité de la fibre
\(T_{\min}^{-1}(U_6)\). Leurs multiplicités satisfont
\[
 432\cdot144+18\cdot432+18\cdot1944=104\,976,
\]
de sorte que la somme de toutes les fibres retrouve exactement le domaine
\(\Omega_{\rm seed}\). La projection directe
\(\rho_3(U_6)\) a \(243\) valeurs distinctes.
\end{theorem}

\paragraph{Démonstration du recensement exhaustif \(104\,976\to468\to243\) et des multiplicités des fibres de \(T_{\min}\)}
La récurrence qui définit \(T_{\min}\) réside dans
\cref{pf84:E02-16} ; la factorisation de son image et le calcul des fibres,
dans \cref{pf84:E02-20,cor:censo-proyeccion-observable}. Cette annexe n'introduit pas
une seconde énumération : elle matérialise, avec tous ses champs, le
catalogue canonique produit par ces résultats.

\section{Catalogue complet des 468 émissions}

Le tableau suivant incorpore les dix champs du catalogue exhaustif.
L'ordinal de la première colonne fixe un ordre de lecture, d'un type distinct
de celui des germes et des identifiants mathématiques. Les colonnes
abrégées conservent
\texttt{diag\_c}, \texttt{hplus\_type}, \texttt{px\_size},
\texttt{px\_res\_type}, \texttt{delta\_diag\_subset},
    \texttt{delta\_orbit}, \texttt{Esig} et \(\rho_3(U_6)\).

\begingroup
\tiny
\setlength{\tabcolsep}{2pt}
\renewcommand{\arraystretch}{1.12}
\begin{longtable}{@{}r
  >{\ttfamily\raggedright\arraybackslash}p{.245\textwidth}
  >{\raggedright\arraybackslash}p{.085\textwidth}
  >{\ttfamily\raggedright\arraybackslash}p{.105\textwidth}
  >{\ttfamily\raggedright\arraybackslash}p{.20\textwidth}
  >{\ttfamily\raggedright\arraybackslash}p{.15\textwidth}@{}}
\caption{Catalogue canonique complet des émissions U6.}
\label{tab:catalogo-emisiones-u6-rev9}\\
\toprule
\# & \normalfont\(U_6\) / \(\rho_3(U_6)\)
   & \normalfont\(\#T_{\min}^{-1}(U_6)\) / \texttt{diag\_c}
   & \normalfont\texttt{hplus\_type} / \texttt{px\_size} / \texttt{px\_res\_type}
   & \normalfont\texttt{delta\_diag\_subset} / \texttt{delta\_orbit}
   & \normalfont\texttt{Esig}\\
\midrule
\endfirsthead
\multicolumn{6}{c}{\tablename\ \thetable\ (suite)}\\
\toprule
\# & \normalfont\(U_6\) / \(\rho_3(U_6)\)
   & \normalfont\(\#T_{\min}^{-1}(U_6)\) / \texttt{diag\_c}
   & \normalfont\texttt{hplus\_type} / \texttt{px\_size} / \texttt{px\_res\_type}
   & \normalfont\texttt{delta\_diag\_subset} / \texttt{delta\_orbit}
   & \normalfont\texttt{Esig}\\
\midrule
\endhead
\midrule
\multicolumn{6}{r}{Suite à la page suivante.}\\
\endfoot
\bottomrule
\endlastfoot
1 & 252|109|252|903|850|850 / 010011 & 144 / 0 & NO / 6 / A & 0,1,4,5 / 0,1,4,5 & 3,9,3,1,7,7 \\
2 & 232|189|292|923|870|810 / 101200 & 144 / 0 & NO / 6 / A & 0,1,4,5 / 0,1,4,5 & 1,7,7,3,9,3 \\
3 & 923|870|810|232|189|292 / 200101 & 144 / 0 & ES / 6 / A & 0,1,4,5 / 0,1,4,5 & 3,9,3,1,7,7 \\
4 & 903|850|850|252|109|252 / 011010 & 144 / 0 & ES / 6 / A & 0,1,4,5 / 0,1,4,5 & 1,7,7,3,9,3 \\
5 & 272|169|272|923|810|870 / 212200 & 144 / 0 & NO / 8 / A & 0,1,4,5 / 0,1,4,5 & 5,5,5,3,3,9 \\
6 & 252|149|212|943|830|830 / 022122 & 144 / 0 & NO / 8 / A & 0,1,4,5 / 0,1,4,5 & 3,3,9,5,5,5 \\
7 & 943|830|830|252|149|212 / 122022 & 144 / 0 & ES / 8 / A & 0,1,4,5 / 0,1,4,5 & 5,5,5,3,3,9 \\
8 & 923|810|870|272|169|272 / 200212 & 144 / 0 & ES / 8 / A & 0,1,4,5 / 0,1,4,5 & 3,3,9,5,5,5 \\
9 & 227|114|227|998|885|885 / 202200 & 1944 / 0 & NO / 45 / H & 0,1,2,3,4,5,6,7,8 / 0,1,2,3,4,5,6,7,8 & 0,0,0,0,0,0 \\
10 & 292|189|232|903|850|850 / 101011 & 144 / 0 & NO / 8 / B & 2,3,6,8 / 0,1,4,6 & 7,7,1,1,7,7 \\
11 & 232|189|292|963|850|890 / 101012 & 144 / 0 & NO / 8 / B & 2,3,6,8 / 0,1,4,6 & 1,7,7,7,7,1 \\
12 & 998|885|885|227|114|227 / 200202 & 1944 / 0 & ES / 45 / H & 0,1,2,3,4,5,6,7,8 / 0,1,2,3,4,5,6,7,8 & 0,0,0,0,0,0 \\
13 & 963|850|890|232|189|292 / 012101 & 144 / 0 & ES / 8 / B & 2,3,6,8 / 0,1,4,6 & 7,7,1,1,7,7 \\
14 & 903|850|850|292|189|232 / 011101 & 144 / 0 & ES / 8 / B & 2,3,6,8 / 0,1,4,6 & 1,7,7,7,7,1 \\
15 & 292|129|292|943|830|830 / 101122 & 144 / 0 & NO / 8 / A & 1,2,3,4 / 0,1,2,3 & 7,1,7,5,5,5 \\
16 & 292|189|232|923|810|870 / 101200 & 144 / 0 & NO / 6 / A & 2,3,7 / 0,1,5 & 7,7,1,3,3,9 \\
17 & 272|169|272|963|890|850 / 212021 & 144 / 0 & NO / 8 / A & 1,2,3,4 / 0,1,2,3 & 5,5,5,7,1,7 \\
18 & 252|149|212|963|850|890 / 022012 & 144 / 0 & NO / 6 / A & 2,3,7 / 0,1,5 & 3,3,9,7,7,1 \\
19 & 963|890|850|272|169|272 / 021212 & 144 / 0 & ES / 8 / A & 1,2,3,4 / 0,1,2,3 & 7,1,7,5,5,5 \\
20 & 963|850|890|252|149|212 / 012022 & 144 / 0 & ES / 6 / A & 2,3,7 / 0,1,5 & 7,7,1,3,3,9 \\
21 & 943|830|830|292|129|292 / 122101 & 144 / 0 & ES / 8 / A & 1,2,3,4 / 0,1,2,3 & 5,5,5,7,1,7 \\
22 & 923|810|870|292|189|232 / 200101 & 144 / 0 & ES / 6 / A & 2,3,7 / 0,1,5 & 3,3,9,7,7,1 \\
23 & 292|189|232|943|830|830 / 101122 & 144 / 0 & NO / 6 / A & 1,4,6,8 / 0,2,4,6 & 7,7,1,5,5,5 \\
24 & 212|149|252|963|890|850 / 220021 & 144 / 0 & NO / 8 / A & 1,4,6,8 / 0,2,4,6 & 9,3,3,7,1,7 \\
25 & 272|169|272|963|850|890 / 212012 & 144 / 0 & NO / 6 / A & 1,4,6,8 / 0,2,4,6 & 5,5,5,7,7,1 \\
26 & 292|129|292|983|810|810 / 101200 & 144 / 0 & NO / 8 / A & 1,4,6,8 / 0,2,4,6 & 7,1,7,9,3,3 \\
27 & 963|850|890|272|169|272 / 012212 & 144 / 0 & ES / 6 / A & 1,4,6,8 / 0,2,4,6 & 7,7,1,5,5,5 \\
28 & 983|810|810|292|129|292 / 200101 & 144 / 0 & ES / 8 / A & 1,4,6,8 / 0,2,4,6 & 9,3,3,7,1,7 \\
29 & 943|830|830|292|189|232 / 122101 & 144 / 0 & ES / 6 / A & 1,4,6,8 / 0,2,4,6 & 5,5,5,7,7,1 \\
30 & 963|890|850|212|149|252 / 021220 & 144 / 0 & ES / 8 / A & 1,4,6,8 / 0,2,4,6 & 7,1,7,9,3,3 \\
31 & 292|129|292|963|890|850 / 101021 & 144 / 0 & NO / 4 / B & 0,5 / 0,4 & 7,1,7,7,1,7 \\
32 & 963|890|850|292|129|292 / 021101 & 144 / 0 & ES / 4 / B & 0,5 / 0,4 & 7,1,7,7,1,7 \\
33 & 272|169|272|983|810|810 / 212200 & 144 / 0 & NO / 6 / A & 6,7,8 / 0,1,2 & 5,5,5,9,3,3 \\
34 & 212|149|252|943|830|830 / 220122 & 144 / 0 & NO / 6 / A & 6,7,8 / 0,1,2 & 9,3,3,5,5,5 \\
35 & 943|830|830|212|149|252 / 122220 & 144 / 0 & ES / 6 / A & 6,7,8 / 0,1,2 & 5,5,5,9,3,3 \\
36 & 983|810|810|272|169|272 / 200212 & 144 / 0 & ES / 6 / A & 6,7,8 / 0,1,2 & 9,3,3,5,5,5 \\
37 & 252|109|252|943|830|830 / 010122 & 144 / 0 & NO / 6 / A & 1,2,3,4 / 0,1,2,3 & 3,9,3,5,5,5 \\
38 & 272|169|272|923|870|810 / 212200 & 144 / 0 & NO / 6 / A & 1,2,3,4 / 0,1,2,3 & 5,5,5,3,9,3 \\
39 & 923|870|810|272|169|272 / 200212 & 144 / 0 & ES / 6 / A & 1,2,3,4 / 0,1,2,3 & 3,9,3,5,5,5 \\
40 & 943|830|830|252|109|252 / 122010 & 144 / 0 & ES / 6 / A & 1,2,3,4 / 0,1,2,3 & 5,5,5,3,9,3 \\
41 & 252|109|252|923|870|810 / 010200 & 144 / 0 & NO / 4 / B & 2,3 / 0,1 & 3,9,3,3,9,3 \\
42 & 923|870|810|252|109|252 / 200010 & 144 / 0 & ES / 4 / B & 2,3 / 0,1 & 3,9,3,3,9,3 \\
43 & 272|169|272|903|850|850 / 212011 & 144 / 0 & NO / 6 / A & 0,5,7 / 0,2,4 & 5,5,5,1,7,7 \\
44 & 232|189|292|943|830|830 / 101122 & 144 / 0 & NO / 6 / A & 0,5,7 / 0,2,4 & 1,7,7,5,5,5 \\
45 & 943|830|830|232|189|292 / 122101 & 144 / 0 & ES / 6 / A & 0,5,7 / 0,2,4 & 5,5,5,1,7,7 \\
46 & 903|850|850|272|169|272 / 011212 & 144 / 0 & ES / 6 / A & 0,5,7 / 0,2,4 & 1,7,7,5,5,5 \\
47 & 212|149|252|923|810|870 / 220200 & 144 / 0 & NO / 8 / B & 0,5,6,8 / 0,1,3,4 & 9,3,3,3,3,9 \\
48 & 252|149|212|983|810|810 / 022200 & 144 / 0 & NO / 8 / B & 0,5,6,8 / 0,1,3,4 & 3,3,9,9,3,3 \\
49 & 983|810|810|252|149|212 / 200022 & 144 / 0 & ES / 8 / B & 0,5,6,8 / 0,1,3,4 & 9,3,3,3,3,9 \\
50 & 923|810|870|212|149|252 / 200220 & 144 / 0 & ES / 8 / B & 0,5,6,8 / 0,1,3,4 & 3,3,9,9,3,3 \\
51 & 272|169|272|943|830|830 / 212122 & 432 / 0 & NO / 12 / B & 0,2,3,5,6,8 / 0,1,3,4,6,7 & 5,5,5,5,5,5 \\
52 & 943|830|830|272|169|272 / 122212 & 432 / 0 & ES / 12 / B & 0,2,3,5,6,8 / 0,1,3,4,6,7 & 5,5,5,5,5,5 \\
53 & 694|438|581|232|189|292 / 102101 & 144 / 1 & NO / 6 / A & 0,3,4,8 / 0,1,4,5 & 3,9,3,1,7,7 \\
54 & 674|418|521|252|109|252 / 212010 & 144 / 1 & NO / 6 / A & 0,3,4,8 / 0,1,4,5 & 1,7,7,3,9,3 \\
55 & 252|109|252|674|418|521 / 010212 & 144 / 1 & ES / 6 / A & 0,3,4,8 / 0,1,4,5 & 3,9,3,1,7,7 \\
56 & 232|189|292|694|438|581 / 101102 & 144 / 1 & ES / 6 / A & 0,3,4,8 / 0,1,4,5 & 1,7,7,3,9,3 \\
57 & 614|498|501|252|149|212 / 200022 & 144 / 1 & NO / 8 / A & 0,3,4,8 / 0,1,4,5 & 5,5,5,3,3,9 \\
58 & 694|478|541|272|169|272 / 111212 & 144 / 1 & NO / 8 / A & 0,3,4,8 / 0,1,4,5 & 3,3,9,5,5,5 \\
59 & 272|169|272|694|478|541 / 212111 & 144 / 1 & ES / 8 / A & 0,3,4,8 / 0,1,4,5 & 5,5,5,3,3,9 \\
60 & 252|149|212|614|498|501 / 022200 & 144 / 1 & ES / 8 / A & 0,3,4,8 / 0,1,4,5 & 3,3,9,5,5,5 \\
61 & 669|443|556|227|114|227 / 021202 & 1944 / 1 & NO / 45 / H & 0,1,2,3,4,5,6,7,8 / 0,1,2,3,4,5,6,7,8 & 0,0,0,0,0,0 \\
62 & 634|418|561|232|189|292 / 110101 & 144 / 1 & NO / 8 / B & 1,2,5,7 / 0,1,4,6 & 7,7,1,1,7,7 \\
63 & 674|418|521|292|189|232 / 212101 & 144 / 1 & NO / 8 / B & 1,2,5,7 / 0,1,4,6 & 1,7,7,7,7,1 \\
64 & 227|114|227|669|443|556 / 202021 & 1944 / 1 & ES / 45 / H & 0,1,2,3,4,5,6,7,8 / 0,1,2,3,4,5,6,7,8 & 0,0,0,0,0,0 \\
65 & 292|189|232|674|418|521 / 101212 & 144 / 1 & ES / 8 / B & 1,2,5,7 / 0,1,4,6 & 7,7,1,1,7,7 \\
66 & 232|189|292|634|418|561 / 101110 & 144 / 1 & ES / 8 / B & 1,2,5,7 / 0,1,4,6 & 1,7,7,7,7,1 \\
67 & 634|458|521|272|169|272 / 122212 & 144 / 1 & NO / 8 / A & 0,1,2,3 / 0,1,2,3 & 7,1,7,5,5,5 \\
68 & 634|418|561|252|149|212 / 110022 & 144 / 1 & NO / 6 / A & 1,2,6 / 0,1,5 & 7,7,1,3,3,9 \\
69 & 614|498|501|292|129|292 / 200101 & 144 / 1 & NO / 8 / A & 0,1,2,3 / 0,1,2,3 & 5,5,5,7,1,7 \\
70 & 694|478|541|292|189|232 / 111101 & 144 / 1 & NO / 6 / A & 1,2,6 / 0,1,5 & 3,3,9,7,7,1 \\
71 & 292|129|292|614|498|501 / 101200 & 144 / 1 & ES / 8 / A & 0,1,2,3 / 0,1,2,3 & 7,1,7,5,5,5 \\
72 & 292|189|232|694|478|541 / 101111 & 144 / 1 & ES / 6 / A & 1,2,6 / 0,1,5 & 7,7,1,3,3,9 \\
73 & 272|169|272|634|458|521 / 212122 & 144 / 1 & ES / 8 / A & 0,1,2,3 / 0,1,2,3 & 5,5,5,7,1,7 \\
74 & 252|149|212|634|418|561 / 022110 & 144 / 1 & ES / 6 / A & 1,2,6 / 0,1,5 & 3,3,9,7,7,1 \\
75 & 634|418|561|272|169|272 / 110212 & 144 / 1 & NO / 6 / A & 0,3,5,7 / 0,2,4,6 & 7,7,1,5,5,5 \\
76 & 654|478|581|292|129|292 / 012101 & 144 / 1 & NO / 8 / A & 0,3,5,7 / 0,2,4,6 & 9,3,3,7,1,7 \\
77 & 614|498|501|292|189|232 / 200101 & 144 / 1 & NO / 6 / A & 0,3,5,7 / 0,2,4,6 & 5,5,5,7,7,1 \\
78 & 634|458|521|212|149|252 / 122220 & 144 / 1 & NO / 8 / A & 0,3,5,7 / 0,2,4,6 & 7,1,7,9,3,3 \\
79 & 292|189|232|614|498|501 / 101200 & 144 / 1 & ES / 6 / A & 0,3,5,7 / 0,2,4,6 & 7,7,1,5,5,5 \\
80 & 212|149|252|634|458|521 / 220122 & 144 / 1 & ES / 8 / A & 0,3,5,7 / 0,2,4,6 & 9,3,3,7,1,7 \\
81 & 272|169|272|634|418|561 / 212110 & 144 / 1 & ES / 6 / A & 0,3,5,7 / 0,2,4,6 & 5,5,5,7,7,1 \\
82 & 292|129|292|654|478|581 / 101012 & 144 / 1 & ES / 8 / A & 0,3,5,7 / 0,2,4,6 & 7,1,7,9,3,3 \\
83 & 634|458|521|292|129|292 / 122101 & 144 / 1 & NO / 4 / B & 4,8 / 0,4 & 7,1,7,7,1,7 \\
84 & 292|129|292|634|458|521 / 101122 & 144 / 1 & ES / 4 / B & 4,8 / 0,4 & 7,1,7,7,1,7 \\
85 & 614|498|501|212|149|252 / 200220 & 144 / 1 & NO / 6 / A & 5,6,7 / 0,1,2 & 5,5,5,9,3,3 \\
86 & 654|478|581|272|169|272 / 012212 & 144 / 1 & NO / 6 / A & 5,6,7 / 0,1,2 & 9,3,3,5,5,5 \\
87 & 272|169|272|654|478|581 / 212012 & 144 / 1 & ES / 6 / A & 5,6,7 / 0,1,2 & 5,5,5,9,3,3 \\
88 & 212|149|252|614|498|501 / 220200 & 144 / 1 & ES / 6 / A & 5,6,7 / 0,1,2 & 9,3,3,5,5,5 \\
89 & 694|438|581|272|169|272 / 102212 & 144 / 1 & NO / 6 / A & 0,1,2,3 / 0,1,2,3 & 3,9,3,5,5,5 \\
90 & 614|498|501|252|109|252 / 200010 & 144 / 1 & NO / 6 / A & 0,1,2,3 / 0,1,2,3 & 5,5,5,3,9,3 \\
91 & 252|109|252|614|498|501 / 010200 & 144 / 1 & ES / 6 / A & 0,1,2,3 / 0,1,2,3 & 3,9,3,5,5,5 \\
92 & 272|169|272|694|438|581 / 212102 & 144 / 1 & ES / 6 / A & 0,1,2,3 / 0,1,2,3 & 5,5,5,3,9,3 \\
93 & 694|438|581|252|109|252 / 102010 & 144 / 1 & NO / 4 / B & 1,2 / 0,1 & 3,9,3,3,9,3 \\
94 & 252|109|252|694|438|581 / 010102 & 144 / 1 & ES / 4 / B & 1,2 / 0,1 & 3,9,3,3,9,3 \\
95 & 614|498|501|232|189|292 / 200101 & 144 / 1 & NO / 6 / A & 4,6,8 / 0,2,4 & 5,5,5,1,7,7 \\
96 & 674|418|521|272|169|272 / 212212 & 144 / 1 & NO / 6 / A & 4,6,8 / 0,2,4 & 1,7,7,5,5,5 \\
97 & 272|169|272|674|418|521 / 212212 & 144 / 1 & ES / 6 / A & 4,6,8 / 0,2,4 & 5,5,5,1,7,7 \\
98 & 232|189|292|614|498|501 / 101200 & 144 / 1 & ES / 6 / A & 4,6,8 / 0,2,4 & 1,7,7,5,5,5 \\
99 & 654|478|581|252|149|212 / 012022 & 144 / 1 & NO / 8 / B & 4,5,7,8 / 0,1,3,4 & 9,3,3,3,3,9 \\
100 & 694|478|541|212|149|252 / 111220 & 144 / 1 & NO / 8 / B & 4,5,7,8 / 0,1,3,4 & 3,3,9,9,3,3 \\
101 & 212|149|252|694|478|541 / 220111 & 144 / 1 & ES / 8 / B & 4,5,7,8 / 0,1,3,4 & 9,3,3,3,3,9 \\
102 & 252|149|212|654|478|581 / 022012 & 144 / 1 & ES / 8 / B & 4,5,7,8 / 0,1,3,4 & 3,3,9,9,3,3 \\
103 & 614|498|501|272|169|272 / 200212 & 432 / 1 & NO / 12 / B & 1,2,4,5,7,8 / 0,1,3,4,6,7 & 5,5,5,5,5,5 \\
104 & 272|169|272|614|498|501 / 212200 & 432 / 1 & ES / 12 / B & 1,2,4,5,7,8 / 0,1,3,4,6,7 & 5,5,5,5,5,5 \\
105 & 923|870|810|561|418|634 / 200011 & 144 / 2 & NO / 6 / A & 2,3,7,8 / 0,1,4,5 & 3,9,3,1,7,7 \\
106 & 903|850|850|581|438|694 / 011201 & 144 / 2 & NO / 6 / A & 2,3,7,8 / 0,1,4,5 & 1,7,7,3,9,3 \\
107 & 581|438|694|903|850|850 / 201011 & 144 / 2 & ES / 6 / A & 2,3,7,8 / 0,1,4,5 & 3,9,3,1,7,7 \\
108 & 561|418|634|923|870|810 / 011200 & 144 / 2 & ES / 6 / A & 2,3,7,8 / 0,1,4,5 & 1,7,7,3,9,3 \\
109 & 943|830|830|581|478|654 / 122210 & 144 / 2 & NO / 8 / A & 2,3,7,8 / 0,1,4,5 & 5,5,5,3,3,9 \\
110 & 923|810|870|501|498|614 / 200002 & 144 / 2 & NO / 8 / A & 2,3,7,8 / 0,1,4,5 & 3,3,9,5,5,5 \\
111 & 501|498|614|923|810|870 / 002200 & 144 / 2 & ES / 8 / A & 2,3,7,8 / 0,1,4,5 & 5,5,5,3,3,9 \\
112 & 581|478|654|943|830|830 / 210122 & 144 / 2 & ES / 8 / A & 2,3,7,8 / 0,1,4,5 & 3,3,9,5,5,5 \\
113 & 998|885|885|556|443|669 / 200120 & 1944 / 2 & NO / 45 / H & 0,1,2,3,4,5,6,7,8 / 0,1,2,3,4,5,6,7,8 & 0,0,0,0,0,0 \\
114 & 963|850|890|561|418|634 / 012011 & 144 / 2 & NO / 8 / B & 0,1,4,6 / 0,1,4,6 & 7,7,1,1,7,7 \\
115 & 903|850|850|521|418|674 / 011212 & 144 / 2 & NO / 8 / B & 0,1,4,6 / 0,1,4,6 & 1,7,7,7,7,1 \\
116 & 556|443|669|998|885|885 / 120200 & 1944 / 2 & ES / 45 / H & 0,1,2,3,4,5,6,7,8 / 0,1,2,3,4,5,6,7,8 & 0,0,0,0,0,0 \\
117 & 521|418|674|903|850|850 / 212011 & 144 / 2 & ES / 8 / B & 0,1,4,6 / 0,1,4,6 & 7,7,1,1,7,7 \\
118 & 561|418|634|963|850|890 / 011012 & 144 / 2 & ES / 8 / B & 0,1,4,6 / 0,1,4,6 & 1,7,7,7,7,1 \\
119 & 963|890|850|501|498|614 / 021002 & 144 / 2 & NO / 8 / A & 0,1,2,8 / 0,1,2,3 & 7,1,7,5,5,5 \\
120 & 963|850|890|581|478|654 / 012210 & 144 / 2 & NO / 6 / A & 0,1,5 / 0,1,5 & 7,7,1,3,3,9 \\
121 & 943|830|830|521|458|634 / 122221 & 144 / 2 & NO / 8 / A & 0,1,2,8 / 0,1,2,3 & 5,5,5,7,1,7 \\
122 & 923|810|870|521|418|674 / 200212 & 144 / 2 & NO / 6 / A & 0,1,5 / 0,1,5 & 3,3,9,7,7,1 \\
123 & 521|458|634|943|830|830 / 221122 & 144 / 2 & ES / 8 / A & 0,1,2,8 / 0,1,2,3 & 7,1,7,5,5,5 \\
124 & 521|418|674|923|810|870 / 212200 & 144 / 2 & ES / 6 / A & 0,1,5 / 0,1,5 & 7,7,1,3,3,9 \\
125 & 501|498|614|963|890|850 / 002021 & 144 / 2 & ES / 8 / A & 0,1,2,8 / 0,1,2,3 & 5,5,5,7,1,7 \\
126 & 581|478|654|963|850|890 / 210012 & 144 / 2 & ES / 6 / A & 0,1,5 / 0,1,5 & 3,3,9,7,7,1 \\
127 & 963|850|890|501|498|614 / 012002 & 144 / 2 & NO / 6 / A & 2,4,6,8 / 0,2,4,6 & 7,7,1,5,5,5 \\
128 & 983|810|810|521|458|634 / 200221 & 144 / 2 & NO / 8 / A & 2,4,6,8 / 0,2,4,6 & 9,3,3,7,1,7 \\
129 & 943|830|830|521|418|674 / 122212 & 144 / 2 & NO / 6 / A & 2,4,6,8 / 0,2,4,6 & 5,5,5,7,7,1 \\
130 & 963|890|850|541|478|694 / 021111 & 144 / 2 & NO / 8 / A & 2,4,6,8 / 0,2,4,6 & 7,1,7,9,3,3 \\
131 & 521|418|674|943|830|830 / 212122 & 144 / 2 & ES / 6 / A & 2,4,6,8 / 0,2,4,6 & 7,7,1,5,5,5 \\
132 & 541|478|694|963|890|850 / 111021 & 144 / 2 & ES / 8 / A & 2,4,6,8 / 0,2,4,6 & 9,3,3,7,1,7 \\
133 & 501|498|614|963|850|890 / 002012 & 144 / 2 & ES / 6 / A & 2,4,6,8 / 0,2,4,6 & 5,5,5,7,7,1 \\
134 & 521|458|634|983|810|810 / 221200 & 144 / 2 & ES / 8 / A & 2,4,6,8 / 0,2,4,6 & 7,1,7,9,3,3 \\
135 & 963|890|850|521|458|634 / 021221 & 144 / 2 & NO / 4 / B & 3,7 / 0,4 & 7,1,7,7,1,7 \\
136 & 521|458|634|963|890|850 / 221021 & 144 / 2 & ES / 4 / B & 3,7 / 0,4 & 7,1,7,7,1,7 \\
137 & 943|830|830|541|478|694 / 122111 & 144 / 2 & NO / 6 / A & 4,5,6 / 0,1,2 & 5,5,5,9,3,3 \\
138 & 983|810|810|501|498|614 / 200002 & 144 / 2 & NO / 6 / A & 4,5,6 / 0,1,2 & 9,3,3,5,5,5 \\
139 & 501|498|614|983|810|810 / 002200 & 144 / 2 & ES / 6 / A & 4,5,6 / 0,1,2 & 5,5,5,9,3,3 \\
140 & 541|478|694|943|830|830 / 111122 & 144 / 2 & ES / 6 / A & 4,5,6 / 0,1,2 & 9,3,3,5,5,5 \\
141 & 923|870|810|501|498|614 / 200002 & 144 / 2 & NO / 6 / A & 0,1,2,8 / 0,1,2,3 & 3,9,3,5,5,5 \\
142 & 943|830|830|581|438|694 / 122201 & 144 / 2 & NO / 6 / A & 0,1,2,8 / 0,1,2,3 & 5,5,5,3,9,3 \\
143 & 581|438|694|943|830|830 / 201122 & 144 / 2 & ES / 6 / A & 0,1,2,8 / 0,1,2,3 & 3,9,3,5,5,5 \\
144 & 501|498|614|923|870|810 / 002200 & 144 / 2 & ES / 6 / A & 0,1,2,8 / 0,1,2,3 & 5,5,5,3,9,3 \\
145 & 923|870|810|581|438|694 / 200201 & 144 / 2 & NO / 4 / B & 0,1 / 0,1 & 3,9,3,3,9,3 \\
146 & 581|438|694|923|870|810 / 201200 & 144 / 2 & ES / 4 / B & 0,1 / 0,1 & 3,9,3,3,9,3 \\
147 & 943|830|830|561|418|634 / 122011 & 144 / 2 & NO / 6 / A & 3,5,7 / 0,2,4 & 5,5,5,1,7,7 \\
148 & 903|850|850|501|498|614 / 011002 & 144 / 2 & NO / 6 / A & 3,5,7 / 0,2,4 & 1,7,7,5,5,5 \\
149 & 501|498|614|903|850|850 / 002011 & 144 / 2 & ES / 6 / A & 3,5,7 / 0,2,4 & 5,5,5,1,7,7 \\
150 & 561|418|634|943|830|830 / 011122 & 144 / 2 & ES / 6 / A & 3,5,7 / 0,2,4 & 1,7,7,5,5,5 \\
151 & 983|810|810|581|478|654 / 200210 & 144 / 2 & NO / 8 / B & 3,4,6,7 / 0,1,3,4 & 9,3,3,3,3,9 \\
152 & 923|810|870|541|478|694 / 200111 & 144 / 2 & NO / 8 / B & 3,4,6,7 / 0,1,3,4 & 3,3,9,9,3,3 \\
153 & 541|478|694|923|810|870 / 111200 & 144 / 2 & ES / 8 / B & 3,4,6,7 / 0,1,3,4 & 9,3,3,3,3,9 \\
154 & 581|478|654|983|810|810 / 210200 & 144 / 2 & ES / 8 / B & 3,4,6,7 / 0,1,3,4 & 3,3,9,9,3,3 \\
155 & 943|830|830|501|498|614 / 122002 & 432 / 2 & NO / 12 / B & 0,1,3,4,6,7 / 0,1,3,4,6,7 & 5,5,5,5,5,5 \\
156 & 501|498|614|943|830|830 / 002122 & 432 / 2 & ES / 12 / B & 0,1,3,4,6,7 / 0,1,3,4,6,7 & 5,5,5,5,5,5 \\
157 & 252|212|149|890|850|963 / 022210 & 144 / 3 & NO / 6 / A & 1,2,6,7 / 0,1,4,5 & 3,9,3,1,7,7 \\
158 & 232|292|189|810|870|923 / 110002 & 144 / 3 & NO / 6 / A & 1,2,6,7 / 0,1,4,5 & 1,7,7,3,9,3 \\
159 & 810|870|923|232|292|189 / 002110 & 144 / 3 & ES / 6 / A & 1,2,6,7 / 0,1,4,5 & 3,9,3,1,7,7 \\
160 & 890|850|963|252|212|149 / 210022 & 144 / 3 & ES / 6 / A & 1,2,6,7 / 0,1,4,5 & 1,7,7,3,9,3 \\
161 & 272|272|169|810|810|983 / 221002 & 144 / 3 & NO / 8 / A & 1,2,6,7 / 0,1,4,5 & 5,5,5,3,3,9 \\
162 & 252|252|109|830|830|943 / 001221 & 144 / 3 & NO / 8 / A & 1,2,6,7 / 0,1,4,5 & 3,3,9,5,5,5 \\
163 & 830|830|943|252|252|109 / 221001 & 144 / 3 & ES / 8 / A & 1,2,6,7 / 0,1,4,5 & 5,5,5,3,3,9 \\
164 & 810|810|983|272|272|169 / 002221 & 144 / 3 & ES / 8 / A & 1,2,6,7 / 0,1,4,5 & 3,3,9,5,5,5 \\
165 & 227|227|114|885|885|998 / 220002 & 1944 / 3 & NO / 45 / H & 0,1,2,3,4,5,6,7,8 / 0,1,2,3,4,5,6,7,8 & 0,0,0,0,0,0 \\
166 & 292|292|129|890|850|963 / 110210 & 144 / 3 & NO / 8 / B & 0,3,5,8 / 0,1,4,6 & 7,7,1,1,7,7 \\
167 & 232|292|189|850|850|903 / 110110 & 144 / 3 & NO / 8 / B & 0,3,5,8 / 0,1,4,6 & 1,7,7,7,7,1 \\
168 & 885|885|998|227|227|114 / 002220 & 1944 / 3 & ES / 45 / H & 0,1,2,3,4,5,6,7,8 / 0,1,2,3,4,5,6,7,8 & 0,0,0,0,0,0 \\
169 & 850|850|903|232|292|189 / 110110 & 144 / 3 & ES / 8 / B & 0,3,5,8 / 0,1,4,6 & 7,7,1,1,7,7 \\
170 & 890|850|963|292|292|129 / 210110 & 144 / 3 & ES / 8 / B & 0,3,5,8 / 0,1,4,6 & 1,7,7,7,7,1 \\
171 & 292|232|189|830|830|943 / 110221 & 144 / 3 & NO / 8 / A & 0,1,7,8 / 0,1,2,3 & 7,1,7,5,5,5 \\
172 & 292|292|129|810|810|983 / 110002 & 144 / 3 & NO / 6 / A & 0,4,8 / 0,1,5 & 7,7,1,3,3,9 \\
173 & 272|272|169|850|890|963 / 221120 & 144 / 3 & NO / 8 / A & 0,1,7,8 / 0,1,2,3 & 5,5,5,7,1,7 \\
174 & 252|252|109|850|850|903 / 001110 & 144 / 3 & NO / 6 / A & 0,4,8 / 0,1,5 & 3,3,9,7,7,1 \\
175 & 850|890|963|272|272|169 / 120221 & 144 / 3 & ES / 8 / A & 0,1,7,8 / 0,1,2,3 & 7,1,7,5,5,5 \\
176 & 850|850|903|252|252|109 / 110001 & 144 / 3 & ES / 6 / A & 0,4,8 / 0,1,5 & 7,7,1,3,3,9 \\
177 & 830|830|943|292|232|189 / 221110 & 144 / 3 & ES / 8 / A & 0,1,7,8 / 0,1,2,3 & 5,5,5,7,1,7 \\
178 & 810|810|983|292|292|129 / 002110 & 144 / 3 & ES / 6 / A & 0,4,8 / 0,1,5 & 3,3,9,7,7,1 \\
179 & 292|292|129|830|830|943 / 110221 & 144 / 3 & NO / 6 / A & 1,3,5,7 / 0,2,4,6 & 7,7,1,5,5,5 \\
180 & 212|252|149|850|890|963 / 202120 & 144 / 3 & NO / 8 / A & 1,3,5,7 / 0,2,4,6 & 9,3,3,7,1,7 \\
181 & 272|272|169|850|850|903 / 221110 & 144 / 3 & NO / 6 / A & 1,3,5,7 / 0,2,4,6 & 5,5,5,7,7,1 \\
182 & 292|232|189|870|810|923 / 110002 & 144 / 3 & NO / 8 / A & 1,3,5,7 / 0,2,4,6 & 7,1,7,9,3,3 \\
183 & 850|850|903|272|272|169 / 110221 & 144 / 3 & ES / 6 / A & 1,3,5,7 / 0,2,4,6 & 7,7,1,5,5,5 \\
184 & 870|810|923|292|232|189 / 002110 & 144 / 3 & ES / 8 / A & 1,3,5,7 / 0,2,4,6 & 9,3,3,7,1,7 \\
185 & 830|830|943|292|292|129 / 221110 & 144 / 3 & ES / 6 / A & 1,3,5,7 / 0,2,4,6 & 5,5,5,7,7,1 \\
186 & 850|890|963|212|252|149 / 120202 & 144 / 3 & ES / 8 / A & 1,3,5,7 / 0,2,4,6 & 7,1,7,9,3,3 \\
187 & 292|232|189|850|890|963 / 110120 & 144 / 3 & NO / 4 / B & 2,6 / 0,4 & 7,1,7,7,1,7 \\
188 & 850|890|963|292|232|189 / 120110 & 144 / 3 & ES / 4 / B & 2,6 / 0,4 & 7,1,7,7,1,7 \\
189 & 272|272|169|870|810|923 / 221002 & 144 / 3 & NO / 6 / A & 3,4,5 / 0,1,2 & 5,5,5,9,3,3 \\
190 & 212|252|149|830|830|943 / 202221 & 144 / 3 & NO / 6 / A & 3,4,5 / 0,1,2 & 9,3,3,5,5,5 \\
191 & 830|830|943|212|252|149 / 221202 & 144 / 3 & ES / 6 / A & 3,4,5 / 0,1,2 & 5,5,5,9,3,3 \\
192 & 870|810|923|272|272|169 / 002221 & 144 / 3 & ES / 6 / A & 3,4,5 / 0,1,2 & 9,3,3,5,5,5 \\
193 & 252|212|149|830|830|943 / 022221 & 144 / 3 & NO / 6 / A & 0,1,7,8 / 0,1,2,3 & 3,9,3,5,5,5 \\
194 & 272|272|169|810|870|923 / 221002 & 144 / 3 & NO / 6 / A & 0,1,7,8 / 0,1,2,3 & 5,5,5,3,9,3 \\
195 & 810|870|923|272|272|169 / 002221 & 144 / 3 & ES / 6 / A & 0,1,7,8 / 0,1,2,3 & 3,9,3,5,5,5 \\
196 & 830|830|943|252|212|149 / 221022 & 144 / 3 & ES / 6 / A & 0,1,7,8 / 0,1,2,3 & 5,5,5,3,9,3 \\
197 & 252|212|149|810|870|923 / 022002 & 144 / 3 & NO / 4 / B & 0,8 / 0,1 & 3,9,3,3,9,3 \\
198 & 810|870|923|252|212|149 / 002022 & 144 / 3 & ES / 4 / B & 0,8 / 0,1 & 3,9,3,3,9,3 \\
199 & 272|272|169|890|850|963 / 221210 & 144 / 3 & NO / 6 / A & 2,4,6 / 0,2,4 & 5,5,5,1,7,7 \\
200 & 232|292|189|830|830|943 / 110221 & 144 / 3 & NO / 6 / A & 2,4,6 / 0,2,4 & 1,7,7,5,5,5 \\
201 & 830|830|943|232|292|189 / 221110 & 144 / 3 & ES / 6 / A & 2,4,6 / 0,2,4 & 5,5,5,1,7,7 \\
202 & 890|850|963|272|272|169 / 210221 & 144 / 3 & ES / 6 / A & 2,4,6 / 0,2,4 & 1,7,7,5,5,5 \\
203 & 212|252|149|810|810|983 / 202002 & 144 / 3 & NO / 8 / B & 2,3,5,6 / 0,1,3,4 & 9,3,3,3,3,9 \\
204 & 252|252|109|870|810|923 / 001002 & 144 / 3 & NO / 8 / B & 2,3,5,6 / 0,1,3,4 & 3,3,9,9,3,3 \\
205 & 870|810|923|252|252|109 / 002001 & 144 / 3 & ES / 8 / B & 2,3,5,6 / 0,1,3,4 & 9,3,3,3,3,9 \\
206 & 810|810|983|212|252|149 / 002202 & 144 / 3 & ES / 8 / B & 2,3,5,6 / 0,1,3,4 & 3,3,9,9,3,3 \\
207 & 272|272|169|830|830|943 / 221221 & 432 / 3 & NO / 12 / B & 0,2,3,5,6,8 / 0,1,3,4,6,7 & 5,5,5,5,5,5 \\
208 & 830|830|943|272|272|169 / 221221 & 432 / 3 & ES / 12 / B & 0,2,3,5,6,8 / 0,1,3,4,6,7 & 5,5,5,5,5,5 \\
209 & 581|654|478|129|292|292 / 201011 & 144 / 4 & NO / 6 / A & 0,1,5,6 / 0,1,4,5 & 3,9,3,1,7,7 \\
210 & 561|634|418|149|212|252 / 011220 & 144 / 4 & NO / 6 / A & 0,1,5,6 / 0,1,4,5 & 1,7,7,3,9,3 \\
211 & 149|212|252|561|634|418 / 220011 & 144 / 4 & ES / 6 / A & 0,1,5,6 / 0,1,4,5 & 3,9,3,1,7,7 \\
212 & 129|292|292|581|654|478 / 011201 & 144 / 4 & ES / 6 / A & 0,1,5,6 / 0,1,4,5 & 1,7,7,3,9,3 \\
213 & 501|614|498|149|252|212 / 020202 & 144 / 4 & NO / 8 / A & 0,1,5,6 / 0,1,4,5 & 5,5,5,3,3,9 \\
214 & 581|694|438|169|272|272 / 210122 & 144 / 4 & NO / 8 / A & 0,1,5,6 / 0,1,4,5 & 3,3,9,5,5,5 \\
215 & 169|272|272|581|694|438 / 122210 & 144 / 4 & ES / 8 / A & 0,1,5,6 / 0,1,4,5 & 5,5,5,3,3,9 \\
216 & 149|252|212|501|614|498 / 202020 & 144 / 4 & ES / 8 / A & 0,1,5,6 / 0,1,4,5 & 3,3,9,5,5,5 \\
217 & 556|669|443|114|227|227 / 102022 & 1944 / 4 & NO / 45 / H & 0,1,2,3,4,5,6,7,8 / 0,1,2,3,4,5,6,7,8 & 0,0,0,0,0,0 \\
218 & 521|634|458|129|292|292 / 212011 & 144 / 4 & NO / 8 / B & 2,4,7,8 / 0,1,4,6 & 7,7,1,1,7,7 \\
219 & 561|634|418|189|292|232 / 011011 & 144 / 4 & NO / 8 / B & 2,4,7,8 / 0,1,4,6 & 1,7,7,7,7,1 \\
220 & 114|227|227|556|669|443 / 022102 & 1944 / 4 & ES / 45 / H & 0,1,2,3,4,5,6,7,8 / 0,1,2,3,4,5,6,7,8 & 0,0,0,0,0,0 \\
221 & 189|292|232|561|634|418 / 011011 & 144 / 4 & ES / 8 / B & 2,4,7,8 / 0,1,4,6 & 7,7,1,1,7,7 \\
222 & 129|292|292|521|634|458 / 011212 & 144 / 4 & ES / 8 / B & 2,4,7,8 / 0,1,4,6 & 1,7,7,7,7,1 \\
223 & 521|674|418|169|272|272 / 221122 & 144 / 4 & NO / 8 / A & 0,6,7,8 / 0,1,2,3 & 7,1,7,5,5,5 \\
224 & 521|634|458|149|252|212 / 212202 & 144 / 4 & NO / 6 / A & 3,7,8 / 0,1,5 & 7,7,1,3,3,9 \\
225 & 501|614|498|189|232|292 / 020011 & 144 / 4 & NO / 8 / A & 0,6,7,8 / 0,1,2,3 & 5,5,5,7,1,7 \\
226 & 581|694|438|189|292|232 / 210011 & 144 / 4 & NO / 6 / A & 3,7,8 / 0,1,5 & 3,3,9,7,7,1 \\
227 & 189|232|292|501|614|498 / 011020 & 144 / 4 & ES / 8 / A & 0,6,7,8 / 0,1,2,3 & 7,1,7,5,5,5 \\
228 & 189|292|232|581|694|438 / 011210 & 144 / 4 & ES / 6 / A & 3,7,8 / 0,1,5 & 7,7,1,3,3,9 \\
229 & 169|272|272|521|674|418 / 122221 & 144 / 4 & ES / 8 / A & 0,6,7,8 / 0,1,2,3 & 5,5,5,7,1,7 \\
230 & 149|252|212|521|634|458 / 202212 & 144 / 4 & ES / 6 / A & 3,7,8 / 0,1,5 & 3,3,9,7,7,1 \\
231 & 521|634|458|169|272|272 / 212122 & 144 / 4 & NO / 6 / A & 0,2,4,6 / 0,2,4,6 & 7,7,1,5,5,5 \\
232 & 541|694|478|189|232|292 / 111011 & 144 / 4 & NO / 8 / A & 0,2,4,6 / 0,2,4,6 & 9,3,3,7,1,7 \\
233 & 501|614|498|189|292|232 / 020011 & 144 / 4 & NO / 6 / A & 0,2,4,6 / 0,2,4,6 & 5,5,5,7,7,1 \\
234 & 521|674|418|109|252|252 / 221100 & 144 / 4 & NO / 8 / A & 0,2,4,6 / 0,2,4,6 & 7,1,7,9,3,3 \\
235 & 189|292|232|501|614|498 / 011020 & 144 / 4 & ES / 6 / A & 0,2,4,6 / 0,2,4,6 & 7,7,1,5,5,5 \\
236 & 109|252|252|521|674|418 / 100221 & 144 / 4 & ES / 8 / A & 0,2,4,6 / 0,2,4,6 & 9,3,3,7,1,7 \\
237 & 169|272|272|521|634|458 / 122212 & 144 / 4 & ES / 6 / A & 0,2,4,6 / 0,2,4,6 & 5,5,5,7,7,1 \\
238 & 189|232|292|541|694|478 / 011111 & 144 / 4 & ES / 8 / A & 0,2,4,6 / 0,2,4,6 & 7,1,7,9,3,3 \\
239 & 521|674|418|189|232|292 / 221011 & 144 / 4 & NO / 4 / B & 1,5 / 0,4 & 7,1,7,7,1,7 \\
240 & 189|232|292|521|674|418 / 011221 & 144 / 4 & ES / 4 / B & 1,5 / 0,4 & 7,1,7,7,1,7 \\
241 & 501|614|498|109|252|252 / 020100 & 144 / 4 & NO / 6 / A & 2,3,4 / 0,1,2 & 5,5,5,9,3,3 \\
242 & 541|694|478|169|272|272 / 111122 & 144 / 4 & NO / 6 / A & 2,3,4 / 0,1,2 & 9,3,3,5,5,5 \\
243 & 169|272|272|541|694|478 / 122111 & 144 / 4 & ES / 6 / A & 2,3,4 / 0,1,2 & 5,5,5,9,3,3 \\
244 & 109|252|252|501|614|498 / 100020 & 144 / 4 & ES / 6 / A & 2,3,4 / 0,1,2 & 9,3,3,5,5,5 \\
245 & 581|654|478|169|272|272 / 201122 & 144 / 4 & NO / 6 / A & 0,6,7,8 / 0,1,2,3 & 3,9,3,5,5,5 \\
246 & 501|614|498|149|212|252 / 020220 & 144 / 4 & NO / 6 / A & 0,6,7,8 / 0,1,2,3 & 5,5,5,3,9,3 \\
247 & 149|212|252|501|614|498 / 220020 & 144 / 4 & ES / 6 / A & 0,6,7,8 / 0,1,2,3 & 3,9,3,5,5,5 \\
248 & 169|272|272|581|654|478 / 122201 & 144 / 4 & ES / 6 / A & 0,6,7,8 / 0,1,2,3 & 5,5,5,3,9,3 \\
249 & 581|654|478|149|212|252 / 201220 & 144 / 4 & NO / 4 / B & 7,8 / 0,1 & 3,9,3,3,9,3 \\
250 & 149|212|252|581|654|478 / 220201 & 144 / 4 & ES / 4 / B & 7,8 / 0,1 & 3,9,3,3,9,3 \\
251 & 501|614|498|129|292|292 / 020011 & 144 / 4 & NO / 6 / A & 1,3,5 / 0,2,4 & 5,5,5,1,7,7 \\
252 & 561|634|418|169|272|272 / 011122 & 144 / 4 & NO / 6 / A & 1,3,5 / 0,2,4 & 1,7,7,5,5,5 \\
253 & 169|272|272|561|634|418 / 122011 & 144 / 4 & ES / 6 / A & 1,3,5 / 0,2,4 & 5,5,5,1,7,7 \\
254 & 129|292|292|501|614|498 / 011020 & 144 / 4 & ES / 6 / A & 1,3,5 / 0,2,4 & 1,7,7,5,5,5 \\
255 & 541|694|478|149|252|212 / 111202 & 144 / 4 & NO / 8 / B & 1,2,4,5 / 0,1,3,4 & 9,3,3,3,3,9 \\
256 & 581|694|438|109|252|252 / 210100 & 144 / 4 & NO / 8 / B & 1,2,4,5 / 0,1,3,4 & 3,3,9,9,3,3 \\
257 & 109|252|252|581|694|438 / 100210 & 144 / 4 & ES / 8 / B & 1,2,4,5 / 0,1,3,4 & 9,3,3,3,3,9 \\
258 & 149|252|212|541|694|478 / 202111 & 144 / 4 & ES / 8 / B & 1,2,4,5 / 0,1,3,4 & 3,3,9,9,3,3 \\
259 & 501|614|498|169|272|272 / 020122 & 432 / 4 & NO / 12 / B & 1,2,4,5,7,8 / 0,1,3,4,6,7 & 5,5,5,5,5,5 \\
260 & 169|272|272|501|614|498 / 122020 & 432 / 4 & ES / 12 / B & 1,2,4,5,7,8 / 0,1,3,4,6,7 & 5,5,5,5,5,5 \\
261 & 810|983|810|458|634|521 / 020212 & 144 / 5 & NO / 6 / A & 0,4,5,8 / 0,1,4,5 & 3,9,3,1,7,7 \\
262 & 890|963|850|478|654|581 / 201102 & 144 / 5 & NO / 6 / A & 0,4,5,8 / 0,1,4,5 & 1,7,7,3,9,3 \\
263 & 478|654|581|890|963|850 / 102201 & 144 / 5 & ES / 6 / A & 0,4,5,8 / 0,1,4,5 & 3,9,3,1,7,7 \\
264 & 458|634|521|810|983|810 / 212020 & 144 / 5 & ES / 6 / A & 0,4,5,8 / 0,1,4,5 & 1,7,7,3,9,3 \\
265 & 830|943|830|478|694|541 / 212111 & 144 / 5 & NO / 8 / A & 0,4,5,8 / 0,1,4,5 & 5,5,5,3,3,9 \\
266 & 810|923|870|498|614|501 / 020020 & 144 / 5 & NO / 8 / A & 0,4,5,8 / 0,1,4,5 & 3,3,9,5,5,5 \\
267 & 498|614|501|810|923|870 / 020020 & 144 / 5 & ES / 8 / A & 0,4,5,8 / 0,1,4,5 & 5,5,5,3,3,9 \\
268 & 478|694|541|830|943|830 / 111212 & 144 / 5 & ES / 8 / A & 0,4,5,8 / 0,1,4,5 & 3,3,9,5,5,5 \\
269 & 885|998|885|443|669|556 / 020201 & 1944 / 5 & NO / 45 / H & 0,1,2,3,4,5,6,7,8 / 0,1,2,3,4,5,6,7,8 & 0,0,0,0,0,0 \\
270 & 850|963|890|458|634|521 / 102212 & 144 / 5 & NO / 8 / B & 1,3,6,7 / 0,1,4,6 & 7,7,1,1,7,7 \\
271 & 890|963|850|418|634|561 / 201110 & 144 / 5 & NO / 8 / B & 1,3,6,7 / 0,1,4,6 & 1,7,7,7,7,1 \\
272 & 443|669|556|885|998|885 / 201020 & 1944 / 5 & ES / 45 / H & 0,1,2,3,4,5,6,7,8 / 0,1,2,3,4,5,6,7,8 & 0,0,0,0,0,0 \\
273 & 418|634|561|890|963|850 / 110201 & 144 / 5 & ES / 8 / B & 1,3,6,7 / 0,1,4,6 & 7,7,1,1,7,7 \\
274 & 458|634|521|850|963|890 / 212102 & 144 / 5 & ES / 8 / B & 1,3,6,7 / 0,1,4,6 & 1,7,7,7,7,1 \\
275 & 850|903|850|498|614|501 / 101020 & 144 / 5 & NO / 8 / A & 5,6,7,8 / 0,1,2,3 & 7,1,7,5,5,5 \\
276 & 850|963|890|478|694|541 / 102111 & 144 / 5 & NO / 6 / A & 2,6,7 / 0,1,5 & 7,7,1,3,3,9 \\
277 & 830|943|830|418|674|521 / 212122 & 144 / 5 & NO / 8 / A & 5,6,7,8 / 0,1,2,3 & 5,5,5,7,1,7 \\
278 & 810|923|870|418|634|561 / 020110 & 144 / 5 & NO / 6 / A & 2,6,7 / 0,1,5 & 3,3,9,7,7,1 \\
279 & 418|674|521|830|943|830 / 122212 & 144 / 5 & ES / 8 / A & 5,6,7,8 / 0,1,2,3 & 7,1,7,5,5,5 \\
280 & 418|634|561|810|923|870 / 110020 & 144 / 5 & ES / 6 / A & 2,6,7 / 0,1,5 & 7,7,1,3,3,9 \\
281 & 498|614|501|850|903|850 / 020101 & 144 / 5 & ES / 8 / A & 5,6,7,8 / 0,1,2,3 & 5,5,5,7,1,7 \\
282 & 478|694|541|850|963|890 / 111102 & 144 / 5 & ES / 6 / A & 2,6,7 / 0,1,5 & 3,3,9,7,7,1 \\
283 & 850|963|890|498|614|501 / 102020 & 144 / 5 & NO / 6 / A & 1,3,5,8 / 0,2,4,6 & 7,7,1,5,5,5 \\
284 & 870|923|810|418|674|521 / 020122 & 144 / 5 & NO / 8 / A & 1,3,5,8 / 0,2,4,6 & 9,3,3,7,1,7 \\
285 & 830|943|830|418|634|561 / 212110 & 144 / 5 & NO / 6 / A & 1,3,5,8 / 0,2,4,6 & 5,5,5,7,7,1 \\
286 & 850|903|850|438|694|581 / 101012 & 144 / 5 & NO / 8 / A & 1,3,5,8 / 0,2,4,6 & 7,1,7,9,3,3 \\
287 & 418|634|561|830|943|830 / 110212 & 144 / 5 & ES / 6 / A & 1,3,5,8 / 0,2,4,6 & 7,7,1,5,5,5 \\
288 & 438|694|581|850|903|850 / 012101 & 144 / 5 & ES / 8 / A & 1,3,5,8 / 0,2,4,6 & 9,3,3,7,1,7 \\
289 & 498|614|501|850|963|890 / 020102 & 144 / 5 & ES / 6 / A & 1,3,5,8 / 0,2,4,6 & 5,5,5,7,7,1 \\
290 & 418|674|521|870|923|810 / 122020 & 144 / 5 & ES / 8 / A & 1,3,5,8 / 0,2,4,6 & 7,1,7,9,3,3 \\
291 & 850|903|850|418|674|521 / 101122 & 144 / 5 & NO / 4 / B & 0,4 / 0,4 & 7,1,7,7,1,7 \\
292 & 418|674|521|850|903|850 / 122101 & 144 / 5 & ES / 4 / B & 0,4 / 0,4 & 7,1,7,7,1,7 \\
293 & 830|943|830|438|694|581 / 212012 & 144 / 5 & NO / 6 / A & 1,2,3 / 0,1,2 & 5,5,5,9,3,3 \\
294 & 870|923|810|498|614|501 / 020020 & 144 / 5 & NO / 6 / A & 1,2,3 / 0,1,2 & 9,3,3,5,5,5 \\
295 & 498|614|501|870|923|810 / 020020 & 144 / 5 & ES / 6 / A & 1,2,3 / 0,1,2 & 5,5,5,9,3,3 \\
296 & 438|694|581|830|943|830 / 012212 & 144 / 5 & ES / 6 / A & 1,2,3 / 0,1,2 & 9,3,3,5,5,5 \\
297 & 810|983|810|498|614|501 / 020020 & 144 / 5 & NO / 6 / A & 5,6,7,8 / 0,1,2,3 & 3,9,3,5,5,5 \\
298 & 830|943|830|478|654|581 / 212102 & 144 / 5 & NO / 6 / A & 5,6,7,8 / 0,1,2,3 & 5,5,5,3,9,3 \\
299 & 478|654|581|830|943|830 / 102212 & 144 / 5 & ES / 6 / A & 5,6,7,8 / 0,1,2,3 & 3,9,3,5,5,5 \\
300 & 498|614|501|810|983|810 / 020020 & 144 / 5 & ES / 6 / A & 5,6,7,8 / 0,1,2,3 & 5,5,5,3,9,3 \\
301 & 810|983|810|478|654|581 / 020102 & 144 / 5 & NO / 4 / B & 6,7 / 0,1 & 3,9,3,3,9,3 \\
302 & 478|654|581|810|983|810 / 102020 & 144 / 5 & ES / 4 / B & 6,7 / 0,1 & 3,9,3,3,9,3 \\
303 & 830|943|830|458|634|521 / 212212 & 144 / 5 & NO / 6 / A & 0,2,4 / 0,2,4 & 5,5,5,1,7,7 \\
304 & 890|963|850|498|614|501 / 201020 & 144 / 5 & NO / 6 / A & 0,2,4 / 0,2,4 & 1,7,7,5,5,5 \\
305 & 498|614|501|890|963|850 / 020201 & 144 / 5 & ES / 6 / A & 0,2,4 / 0,2,4 & 5,5,5,1,7,7 \\
306 & 458|634|521|830|943|830 / 212212 & 144 / 5 & ES / 6 / A & 0,2,4 / 0,2,4 & 1,7,7,5,5,5 \\
307 & 870|923|810|478|694|541 / 020111 & 144 / 5 & NO / 8 / B & 0,1,3,4 / 0,1,3,4 & 9,3,3,3,3,9 \\
308 & 810|923|870|438|694|581 / 020012 & 144 / 5 & NO / 8 / B & 0,1,3,4 / 0,1,3,4 & 3,3,9,9,3,3 \\
309 & 438|694|581|810|923|870 / 012020 & 144 / 5 & ES / 8 / B & 0,1,3,4 / 0,1,3,4 & 9,3,3,3,3,9 \\
310 & 478|694|541|870|923|810 / 111020 & 144 / 5 & ES / 8 / B & 0,1,3,4 / 0,1,3,4 & 3,3,9,9,3,3 \\
311 & 830|943|830|498|614|501 / 212020 & 432 / 5 & NO / 12 / B & 0,1,3,4,6,7 / 0,1,3,4,6,7 & 5,5,5,5,5,5 \\
312 & 498|614|501|830|943|830 / 020212 & 432 / 5 & ES / 12 / B & 0,1,3,4,6,7 / 0,1,3,4,6,7 & 5,5,5,5,5,5 \\
313 & 149|212|252|890|963|850 / 220201 & 144 / 6 & NO / 6 / A & 3,4,7,8 / 0,1,4,5 & 3,9,3,1,7,7 \\
314 & 129|292|292|810|983|810 / 011020 & 144 / 6 & NO / 6 / A & 3,4,7,8 / 0,1,4,5 & 1,7,7,3,9,3 \\
315 & 810|983|810|129|292|292 / 020011 & 144 / 6 & ES / 6 / A & 3,4,7,8 / 0,1,4,5 & 3,9,3,1,7,7 \\
316 & 890|963|850|149|212|252 / 201220 & 144 / 6 & ES / 6 / A & 3,4,7,8 / 0,1,4,5 & 1,7,7,3,9,3 \\
317 & 169|272|272|810|923|870 / 122020 & 144 / 6 & NO / 8 / A & 3,4,7,8 / 0,1,4,5 & 5,5,5,3,3,9 \\
318 & 149|252|212|830|943|830 / 202212 & 144 / 6 & NO / 8 / A & 3,4,7,8 / 0,1,4,5 & 3,3,9,5,5,5 \\
319 & 830|943|830|149|252|212 / 212202 & 144 / 6 & ES / 8 / A & 3,4,7,8 / 0,1,4,5 & 5,5,5,3,3,9 \\
320 & 810|923|870|169|272|272 / 020122 & 144 / 6 & ES / 8 / A & 3,4,7,8 / 0,1,4,5 & 3,3,9,5,5,5 \\
321 & 114|227|227|885|998|885 / 022020 & 1944 / 6 & NO / 45 / H & 0,1,2,3,4,5,6,7,8 / 0,1,2,3,4,5,6,7,8 & 0,0,0,0,0,0 \\
322 & 189|292|232|890|963|850 / 011201 & 144 / 6 & NO / 8 / B & 0,2,5,6 / 0,1,4,6 & 7,7,1,1,7,7 \\
323 & 129|292|292|850|963|890 / 011102 & 144 / 6 & NO / 8 / B & 0,2,5,6 / 0,1,4,6 & 1,7,7,7,7,1 \\
324 & 885|998|885|114|227|227 / 020022 & 1944 / 6 & ES / 45 / H & 0,1,2,3,4,5,6,7,8 / 0,1,2,3,4,5,6,7,8 & 0,0,0,0,0,0 \\
325 & 850|963|890|129|292|292 / 102011 & 144 / 6 & ES / 8 / B & 0,2,5,6 / 0,1,4,6 & 7,7,1,1,7,7 \\
326 & 890|963|850|189|292|232 / 201011 & 144 / 6 & ES / 8 / B & 0,2,5,6 / 0,1,4,6 & 1,7,7,7,7,1 \\
327 & 189|232|292|830|943|830 / 011212 & 144 / 6 & NO / 8 / A & 4,5,6,7 / 0,1,2,3 & 7,1,7,5,5,5 \\
328 & 189|292|232|810|923|870 / 011020 & 144 / 6 & NO / 6 / A & 1,5,6 / 0,1,5 & 7,7,1,3,3,9 \\
329 & 169|272|272|850|903|850 / 122101 & 144 / 6 & NO / 8 / A & 4,5,6,7 / 0,1,2,3 & 5,5,5,7,1,7 \\
330 & 149|252|212|850|963|890 / 202102 & 144 / 6 & NO / 6 / A & 1,5,6 / 0,1,5 & 3,3,9,7,7,1 \\
331 & 850|903|850|169|272|272 / 101122 & 144 / 6 & ES / 8 / A & 4,5,6,7 / 0,1,2,3 & 7,1,7,5,5,5 \\
332 & 850|963|890|149|252|212 / 102202 & 144 / 6 & ES / 6 / A & 1,5,6 / 0,1,5 & 7,7,1,3,3,9 \\
333 & 830|943|830|189|232|292 / 212011 & 144 / 6 & ES / 8 / A & 4,5,6,7 / 0,1,2,3 & 5,5,5,7,1,7 \\
334 & 810|923|870|189|292|232 / 020011 & 144 / 6 & ES / 6 / A & 1,5,6 / 0,1,5 & 3,3,9,7,7,1 \\
335 & 189|292|232|830|943|830 / 011212 & 144 / 6 & NO / 6 / A & 0,2,4,7 / 0,2,4,6 & 7,7,1,5,5,5 \\
336 & 109|252|252|850|903|850 / 100101 & 144 / 6 & NO / 8 / A & 0,2,4,7 / 0,2,4,6 & 9,3,3,7,1,7 \\
337 & 169|272|272|850|963|890 / 122102 & 144 / 6 & NO / 6 / A & 0,2,4,7 / 0,2,4,6 & 5,5,5,7,7,1 \\
338 & 189|232|292|870|923|810 / 011020 & 144 / 6 & NO / 8 / A & 0,2,4,7 / 0,2,4,6 & 7,1,7,9,3,3 \\
339 & 850|963|890|169|272|272 / 102122 & 144 / 6 & ES / 6 / A & 0,2,4,7 / 0,2,4,6 & 7,7,1,5,5,5 \\
340 & 870|923|810|189|232|292 / 020011 & 144 / 6 & ES / 8 / A & 0,2,4,7 / 0,2,4,6 & 9,3,3,7,1,7 \\
341 & 830|943|830|189|292|232 / 212011 & 144 / 6 & ES / 6 / A & 0,2,4,7 / 0,2,4,6 & 5,5,5,7,7,1 \\
342 & 850|903|850|109|252|252 / 101100 & 144 / 6 & ES / 8 / A & 0,2,4,7 / 0,2,4,6 & 7,1,7,9,3,3 \\
343 & 189|232|292|850|903|850 / 011101 & 144 / 6 & NO / 4 / B & 3,8 / 0,4 & 7,1,7,7,1,7 \\
344 & 850|903|850|189|232|292 / 101011 & 144 / 6 & ES / 4 / B & 3,8 / 0,4 & 7,1,7,7,1,7 \\
345 & 169|272|272|870|923|810 / 122020 & 144 / 6 & NO / 6 / A & 0,1,2 / 0,1,2 & 5,5,5,9,3,3 \\
346 & 109|252|252|830|943|830 / 100212 & 144 / 6 & NO / 6 / A & 0,1,2 / 0,1,2 & 9,3,3,5,5,5 \\
347 & 830|943|830|109|252|252 / 212100 & 144 / 6 & ES / 6 / A & 0,1,2 / 0,1,2 & 5,5,5,9,3,3 \\
348 & 870|923|810|169|272|272 / 020122 & 144 / 6 & ES / 6 / A & 0,1,2 / 0,1,2 & 9,3,3,5,5,5 \\
349 & 149|212|252|830|943|830 / 220212 & 144 / 6 & NO / 6 / A & 4,5,6,7 / 0,1,2,3 & 3,9,3,5,5,5 \\
350 & 169|272|272|810|983|810 / 122020 & 144 / 6 & NO / 6 / A & 4,5,6,7 / 0,1,2,3 & 5,5,5,3,9,3 \\
351 & 810|983|810|169|272|272 / 020122 & 144 / 6 & ES / 6 / A & 4,5,6,7 / 0,1,2,3 & 3,9,3,5,5,5 \\
352 & 830|943|830|149|212|252 / 212220 & 144 / 6 & ES / 6 / A & 4,5,6,7 / 0,1,2,3 & 5,5,5,3,9,3 \\
353 & 149|212|252|810|983|810 / 220020 & 144 / 6 & NO / 4 / B & 5,6 / 0,1 & 3,9,3,3,9,3 \\
354 & 810|983|810|149|212|252 / 020220 & 144 / 6 & ES / 4 / B & 5,6 / 0,1 & 3,9,3,3,9,3 \\
355 & 169|272|272|890|963|850 / 122201 & 144 / 6 & NO / 6 / A & 1,3,8 / 0,2,4 & 5,5,5,1,7,7 \\
356 & 129|292|292|830|943|830 / 011212 & 144 / 6 & NO / 6 / A & 1,3,8 / 0,2,4 & 1,7,7,5,5,5 \\
357 & 830|943|830|129|292|292 / 212011 & 144 / 6 & ES / 6 / A & 1,3,8 / 0,2,4 & 5,5,5,1,7,7 \\
358 & 890|963|850|169|272|272 / 201122 & 144 / 6 & ES / 6 / A & 1,3,8 / 0,2,4 & 1,7,7,5,5,5 \\
359 & 109|252|252|810|923|870 / 100020 & 144 / 6 & NO / 8 / B & 0,2,3,8 / 0,1,3,4 & 9,3,3,3,3,9 \\
360 & 149|252|212|870|923|810 / 202020 & 144 / 6 & NO / 8 / B & 0,2,3,8 / 0,1,3,4 & 3,3,9,9,3,3 \\
361 & 870|923|810|149|252|212 / 020202 & 144 / 6 & ES / 8 / B & 0,2,3,8 / 0,1,3,4 & 9,3,3,3,3,9 \\
362 & 810|923|870|109|252|252 / 020100 & 144 / 6 & ES / 8 / B & 0,2,3,8 / 0,1,3,4 & 3,3,9,9,3,3 \\
363 & 169|272|272|830|943|830 / 122212 & 432 / 6 & NO / 12 / B & 0,2,3,5,6,8 / 0,1,3,4,6,7 & 5,5,5,5,5,5 \\
364 & 830|943|830|169|272|272 / 212122 & 432 / 6 & ES / 12 / B & 0,2,3,5,6,8 / 0,1,3,4,6,7 & 5,5,5,5,5,5 \\
365 & 478|541|694|232|292|189 / 111110 & 144 / 7 & NO / 6 / A & 2,3,6,7 / 0,1,4,5 & 3,9,3,1,7,7 \\
366 & 458|521|634|252|212|149 / 221022 & 144 / 7 & NO / 6 / A & 2,3,6,7 / 0,1,4,5 & 1,7,7,3,9,3 \\
367 & 252|212|149|458|521|634 / 022221 & 144 / 7 & ES / 6 / A & 2,3,6,7 / 0,1,4,5 & 3,9,3,1,7,7 \\
368 & 232|292|189|478|541|694 / 110111 & 144 / 7 & ES / 6 / A & 2,3,6,7 / 0,1,4,5 & 1,7,7,3,9,3 \\
369 & 498|501|614|252|252|109 / 002001 & 144 / 7 & NO / 8 / A & 2,3,6,7 / 0,1,4,5 & 5,5,5,3,3,9 \\
370 & 478|581|654|272|272|169 / 120221 & 144 / 7 & NO / 8 / A & 2,3,6,7 / 0,1,4,5 & 3,3,9,5,5,5 \\
371 & 272|272|169|478|581|654 / 221120 & 144 / 7 & ES / 8 / A & 2,3,6,7 / 0,1,4,5 & 5,5,5,3,3,9 \\
372 & 252|252|109|498|501|614 / 001002 & 144 / 7 & ES / 8 / A & 2,3,6,7 / 0,1,4,5 & 3,3,9,5,5,5 \\
373 & 443|556|669|227|227|114 / 210220 & 1944 / 7 & NO / 45 / H & 0,1,2,3,4,5,6,7,8 / 0,1,2,3,4,5,6,7,8 & 0,0,0,0,0,0 \\
374 & 418|521|674|232|292|189 / 122110 & 144 / 7 & NO / 8 / B & 1,4,5,8 / 0,1,4,6 & 7,7,1,1,7,7 \\
375 & 458|521|634|292|292|129 / 221110 & 144 / 7 & NO / 8 / B & 1,4,5,8 / 0,1,4,6 & 1,7,7,7,7,1 \\
376 & 227|227|114|443|556|669 / 220210 & 1944 / 7 & ES / 45 / H & 0,1,2,3,4,5,6,7,8 / 0,1,2,3,4,5,6,7,8 & 0,0,0,0,0,0 \\
377 & 292|292|129|458|521|634 / 110221 & 144 / 7 & ES / 8 / B & 1,4,5,8 / 0,1,4,6 & 7,7,1,1,7,7 \\
378 & 232|292|189|418|521|674 / 110122 & 144 / 7 & ES / 8 / B & 1,4,5,8 / 0,1,4,6 & 1,7,7,7,7,1 \\
379 & 418|561|634|272|272|169 / 101221 & 144 / 7 & NO / 8 / A & 3,4,5,6 / 0,1,2,3 & 7,1,7,5,5,5 \\
380 & 418|521|674|252|252|109 / 122001 & 144 / 7 & NO / 6 / A & 0,4,5 / 0,1,5 & 7,7,1,3,3,9 \\
381 & 498|501|614|292|232|189 / 002110 & 144 / 7 & NO / 8 / A & 3,4,5,6 / 0,1,2,3 & 5,5,5,7,1,7 \\
382 & 478|581|654|292|292|129 / 120110 & 144 / 7 & NO / 6 / A & 0,4,5 / 0,1,5 & 3,3,9,7,7,1 \\
383 & 292|232|189|498|501|614 / 110002 & 144 / 7 & ES / 8 / A & 3,4,5,6 / 0,1,2,3 & 7,1,7,5,5,5 \\
384 & 292|292|129|478|581|654 / 110120 & 144 / 7 & ES / 6 / A & 0,4,5 / 0,1,5 & 7,7,1,3,3,9 \\
385 & 272|272|169|418|561|634 / 221101 & 144 / 7 & ES / 8 / A & 3,4,5,6 / 0,1,2,3 & 5,5,5,7,1,7 \\
386 & 252|252|109|418|521|674 / 001122 & 144 / 7 & ES / 6 / A & 0,4,5 / 0,1,5 & 3,3,9,7,7,1 \\
387 & 418|521|674|272|272|169 / 122221 & 144 / 7 & NO / 6 / A & 1,3,6,8 / 0,2,4,6 & 7,7,1,5,5,5 \\
388 & 438|581|694|292|232|189 / 021110 & 144 / 7 & NO / 8 / A & 1,3,6,8 / 0,2,4,6 & 9,3,3,7,1,7 \\
389 & 498|501|614|292|292|129 / 002110 & 144 / 7 & NO / 6 / A & 1,3,6,8 / 0,2,4,6 & 5,5,5,7,7,1 \\
390 & 418|561|634|212|252|149 / 101202 & 144 / 7 & NO / 8 / A & 1,3,6,8 / 0,2,4,6 & 7,1,7,9,3,3 \\
391 & 292|292|129|498|501|614 / 110002 & 144 / 7 & ES / 6 / A & 1,3,6,8 / 0,2,4,6 & 7,7,1,5,5,5 \\
392 & 212|252|149|418|561|634 / 202101 & 144 / 7 & ES / 8 / A & 1,3,6,8 / 0,2,4,6 & 9,3,3,7,1,7 \\
393 & 272|272|169|418|521|674 / 221122 & 144 / 7 & ES / 6 / A & 1,3,6,8 / 0,2,4,6 & 5,5,5,7,7,1 \\
394 & 292|232|189|438|581|694 / 110021 & 144 / 7 & ES / 8 / A & 1,3,6,8 / 0,2,4,6 & 7,1,7,9,3,3 \\
395 & 418|561|634|292|232|189 / 101110 & 144 / 7 & NO / 4 / B & 2,7 / 0,4 & 7,1,7,7,1,7 \\
396 & 292|232|189|418|561|634 / 110101 & 144 / 7 & ES / 4 / B & 2,7 / 0,4 & 7,1,7,7,1,7 \\
397 & 498|501|614|212|252|149 / 002202 & 144 / 7 & NO / 6 / A & 0,1,8 / 0,1,2 & 5,5,5,9,3,3 \\
398 & 438|581|694|272|272|169 / 021221 & 144 / 7 & NO / 6 / A & 0,1,8 / 0,1,2 & 9,3,3,5,5,5 \\
399 & 272|272|169|438|581|694 / 221021 & 144 / 7 & ES / 6 / A & 0,1,8 / 0,1,2 & 5,5,5,9,3,3 \\
400 & 212|252|149|498|501|614 / 202002 & 144 / 7 & ES / 6 / A & 0,1,8 / 0,1,2 & 9,3,3,5,5,5 \\
401 & 478|541|694|272|272|169 / 111221 & 144 / 7 & NO / 6 / A & 3,4,5,6 / 0,1,2,3 & 3,9,3,5,5,5 \\
402 & 498|501|614|252|212|149 / 002022 & 144 / 7 & NO / 6 / A & 3,4,5,6 / 0,1,2,3 & 5,5,5,3,9,3 \\
403 & 252|212|149|498|501|614 / 022002 & 144 / 7 & ES / 6 / A & 3,4,5,6 / 0,1,2,3 & 3,9,3,5,5,5 \\
404 & 272|272|169|478|541|694 / 221111 & 144 / 7 & ES / 6 / A & 3,4,5,6 / 0,1,2,3 & 5,5,5,3,9,3 \\
405 & 478|541|694|252|212|149 / 111022 & 144 / 7 & NO / 4 / B & 4,5 / 0,1 & 3,9,3,3,9,3 \\
406 & 252|212|149|478|541|694 / 022111 & 144 / 7 & ES / 4 / B & 4,5 / 0,1 & 3,9,3,3,9,3 \\
407 & 498|501|614|232|292|189 / 002110 & 144 / 7 & NO / 6 / A & 0,2,7 / 0,2,4 & 5,5,5,1,7,7 \\
408 & 458|521|634|272|272|169 / 221221 & 144 / 7 & NO / 6 / A & 0,2,7 / 0,2,4 & 1,7,7,5,5,5 \\
409 & 272|272|169|458|521|634 / 221221 & 144 / 7 & ES / 6 / A & 0,2,7 / 0,2,4 & 5,5,5,1,7,7 \\
410 & 232|292|189|498|501|614 / 110002 & 144 / 7 & ES / 6 / A & 0,2,7 / 0,2,4 & 1,7,7,5,5,5 \\
411 & 438|581|694|252|252|109 / 021001 & 144 / 7 & NO / 8 / B & 1,2,7,8 / 0,1,3,4 & 9,3,3,3,3,9 \\
412 & 478|581|654|212|252|149 / 120202 & 144 / 7 & NO / 8 / B & 1,2,7,8 / 0,1,3,4 & 3,3,9,9,3,3 \\
413 & 212|252|149|478|581|654 / 202120 & 144 / 7 & ES / 8 / B & 1,2,7,8 / 0,1,3,4 & 9,3,3,3,3,9 \\
414 & 252|252|109|438|581|694 / 001021 & 144 / 7 & ES / 8 / B & 1,2,7,8 / 0,1,3,4 & 3,3,9,9,3,3 \\
415 & 498|501|614|272|272|169 / 002221 & 432 / 7 & NO / 12 / B & 1,2,4,5,7,8 / 0,1,3,4,6,7 & 5,5,5,5,5,5 \\
416 & 272|272|169|498|501|614 / 221002 & 432 / 7 & ES / 12 / B & 1,2,4,5,7,8 / 0,1,3,4,6,7 & 5,5,5,5,5,5 \\
417 & 810|870|923|674|521|418 / 002221 & 144 / 8 & NO / 6 / A & 1,2,5,6 / 0,1,4,5 & 3,9,3,1,7,7 \\
418 & 890|850|963|694|541|478 / 210111 & 144 / 8 & NO / 6 / A & 1,2,5,6 / 0,1,4,5 & 1,7,7,3,9,3 \\
419 & 694|541|478|890|850|963 / 111210 & 144 / 8 & ES / 6 / A & 1,2,5,6 / 0,1,4,5 & 3,9,3,1,7,7 \\
420 & 674|521|418|810|870|923 / 221002 & 144 / 8 & ES / 6 / A & 1,2,5,6 / 0,1,4,5 & 1,7,7,3,9,3 \\
421 & 830|830|943|694|581|438 / 221120 & 144 / 8 & NO / 8 / A & 1,2,5,6 / 0,1,4,5 & 5,5,5,3,3,9 \\
422 & 810|810|983|614|501|498 / 002200 & 144 / 8 & NO / 8 / A & 1,2,5,6 / 0,1,4,5 & 3,3,9,5,5,5 \\
423 & 614|501|498|810|810|983 / 200002 & 144 / 8 & ES / 8 / A & 1,2,5,6 / 0,1,4,5 & 5,5,5,3,3,9 \\
424 & 694|581|438|830|830|943 / 120221 & 144 / 8 & ES / 8 / A & 1,2,5,6 / 0,1,4,5 & 3,3,9,5,5,5 \\
425 & 885|885|998|669|556|443 / 002012 & 1944 / 8 & NO / 45 / H & 0,1,2,3,4,5,6,7,8 / 0,1,2,3,4,5,6,7,8 & 0,0,0,0,0,0 \\
426 & 850|850|903|674|521|418 / 110221 & 144 / 8 & NO / 8 / B & 0,3,4,7 / 0,1,4,6 & 7,7,1,1,7,7 \\
427 & 890|850|963|634|521|458 / 210122 & 144 / 8 & NO / 8 / B & 0,3,4,7 / 0,1,4,6 & 1,7,7,7,7,1 \\
428 & 669|556|443|885|885|998 / 012002 & 1944 / 8 & ES / 45 / H & 0,1,2,3,4,5,6,7,8 / 0,1,2,3,4,5,6,7,8 & 0,0,0,0,0,0 \\
429 & 634|521|458|890|850|963 / 122210 & 144 / 8 & ES / 8 / B & 0,3,4,7 / 0,1,4,6 & 7,7,1,1,7,7 \\
430 & 674|521|418|850|850|903 / 221110 & 144 / 8 & ES / 8 / B & 0,3,4,7 / 0,1,4,6 & 1,7,7,7,7,1 \\
431 & 850|890|963|614|501|498 / 120200 & 144 / 8 & NO / 8 / A & 2,3,4,5 / 0,1,2,3 & 7,1,7,5,5,5 \\
432 & 850|850|903|694|581|438 / 110120 & 144 / 8 & NO / 6 / A & 3,4,8 / 0,1,5 & 7,7,1,3,3,9 \\
433 & 830|830|943|634|561|418 / 221101 & 144 / 8 & NO / 8 / A & 2,3,4,5 / 0,1,2,3 & 5,5,5,7,1,7 \\
434 & 810|810|983|634|521|458 / 002122 & 144 / 8 & NO / 6 / A & 3,4,8 / 0,1,5 & 3,3,9,7,7,1 \\
435 & 634|561|418|830|830|943 / 101221 & 144 / 8 & ES / 8 / A & 2,3,4,5 / 0,1,2,3 & 7,1,7,5,5,5 \\
436 & 634|521|458|810|810|983 / 122002 & 144 / 8 & ES / 6 / A & 3,4,8 / 0,1,5 & 7,7,1,3,3,9 \\
437 & 614|501|498|850|890|963 / 200120 & 144 / 8 & ES / 8 / A & 2,3,4,5 / 0,1,2,3 & 5,5,5,7,1,7 \\
438 & 694|581|438|850|850|903 / 120110 & 144 / 8 & ES / 6 / A & 3,4,8 / 0,1,5 & 3,3,9,7,7,1 \\
439 & 850|850|903|614|501|498 / 110200 & 144 / 8 & NO / 6 / A & 0,2,5,7 / 0,2,4,6 & 7,7,1,5,5,5 \\
440 & 870|810|923|634|561|418 / 002101 & 144 / 8 & NO / 8 / A & 0,2,5,7 / 0,2,4,6 & 9,3,3,7,1,7 \\
441 & 830|830|943|634|521|458 / 221122 & 144 / 8 & NO / 6 / A & 0,2,5,7 / 0,2,4,6 & 5,5,5,7,7,1 \\
442 & 850|890|963|654|581|478 / 120021 & 144 / 8 & NO / 8 / A & 0,2,5,7 / 0,2,4,6 & 7,1,7,9,3,3 \\
443 & 634|521|458|830|830|943 / 122221 & 144 / 8 & ES / 6 / A & 0,2,5,7 / 0,2,4,6 & 7,7,1,5,5,5 \\
444 & 654|581|478|850|890|963 / 021120 & 144 / 8 & ES / 8 / A & 0,2,5,7 / 0,2,4,6 & 9,3,3,7,1,7 \\
445 & 614|501|498|850|850|903 / 200110 & 144 / 8 & ES / 6 / A & 0,2,5,7 / 0,2,4,6 & 5,5,5,7,7,1 \\
446 & 634|561|418|870|810|923 / 101002 & 144 / 8 & ES / 8 / A & 0,2,5,7 / 0,2,4,6 & 7,1,7,9,3,3 \\
447 & 850|890|963|634|561|418 / 120101 & 144 / 8 & NO / 4 / B & 1,6 / 0,4 & 7,1,7,7,1,7 \\
448 & 634|561|418|850|890|963 / 101120 & 144 / 8 & ES / 4 / B & 1,6 / 0,4 & 7,1,7,7,1,7 \\
449 & 830|830|943|654|581|478 / 221021 & 144 / 8 & NO / 6 / A & 0,7,8 / 0,1,2 & 5,5,5,9,3,3 \\
450 & 870|810|923|614|501|498 / 002200 & 144 / 8 & NO / 6 / A & 0,7,8 / 0,1,2 & 9,3,3,5,5,5 \\
451 & 614|501|498|870|810|923 / 200002 & 144 / 8 & ES / 6 / A & 0,7,8 / 0,1,2 & 5,5,5,9,3,3 \\
452 & 654|581|478|830|830|943 / 021221 & 144 / 8 & ES / 6 / A & 0,7,8 / 0,1,2 & 9,3,3,5,5,5 \\
453 & 810|870|923|614|501|498 / 002200 & 144 / 8 & NO / 6 / A & 2,3,4,5 / 0,1,2,3 & 3,9,3,5,5,5 \\
454 & 830|830|943|694|541|478 / 221111 & 144 / 8 & NO / 6 / A & 2,3,4,5 / 0,1,2,3 & 5,5,5,3,9,3 \\
455 & 694|541|478|830|830|943 / 111221 & 144 / 8 & ES / 6 / A & 2,3,4,5 / 0,1,2,3 & 3,9,3,5,5,5 \\
456 & 614|501|498|810|870|923 / 200002 & 144 / 8 & ES / 6 / A & 2,3,4,5 / 0,1,2,3 & 5,5,5,3,9,3 \\
457 & 810|870|923|694|541|478 / 002111 & 144 / 8 & NO / 4 / B & 3,4 / 0,1 & 3,9,3,3,9,3 \\
458 & 694|541|478|810|870|923 / 111002 & 144 / 8 & ES / 4 / B & 3,4 / 0,1 & 3,9,3,3,9,3 \\
459 & 830|830|943|674|521|418 / 221221 & 144 / 8 & NO / 6 / A & 1,6,8 / 0,2,4 & 5,5,5,1,7,7 \\
460 & 890|850|963|614|501|498 / 210200 & 144 / 8 & NO / 6 / A & 1,6,8 / 0,2,4 & 1,7,7,5,5,5 \\
461 & 614|501|498|890|850|963 / 200210 & 144 / 8 & ES / 6 / A & 1,6,8 / 0,2,4 & 5,5,5,1,7,7 \\
462 & 674|521|418|830|830|943 / 221221 & 144 / 8 & ES / 6 / A & 1,6,8 / 0,2,4 & 1,7,7,5,5,5 \\
463 & 870|810|923|694|581|438 / 002120 & 144 / 8 & NO / 8 / B & 0,1,6,7 / 0,1,3,4 & 9,3,3,3,3,9 \\
464 & 810|810|983|654|581|478 / 002021 & 144 / 8 & NO / 8 / B & 0,1,6,7 / 0,1,3,4 & 3,3,9,9,3,3 \\
465 & 654|581|478|810|810|983 / 021002 & 144 / 8 & ES / 8 / B & 0,1,6,7 / 0,1,3,4 & 9,3,3,3,3,9 \\
466 & 694|581|438|870|810|923 / 120002 & 144 / 8 & ES / 8 / B & 0,1,6,7 / 0,1,3,4 & 3,3,9,9,3,3 \\
467 & 830|830|943|614|501|498 / 221200 & 432 / 8 & NO / 12 / B & 0,1,3,4,6,7 / 0,1,3,4,6,7 & 5,5,5,5,5,5 \\
468 & 614|501|498|830|830|943 / 200221 & 432 / 8 & ES / 12 / B & 0,1,3,4,6,7 / 0,1,3,4,6,7 & 5,5,5,5,5,5 \\
\end{longtable}
\endgroup

Le tableau est une énumération de l'image de \(T_{\min}\). La spécification
du domaine demeure dans \(\Omega_{\rm seed}=\mathcal S^2\) ; pour retrouver
les conditions initiales individuelles d'une émission, on prend sa fibre
\(T_{\min}^{-1}(U_6)\), dont la cardinalité figure sous
\(\#T_{\min}^{-1}(U_6)\).

\end{HMTScientificOwnerSubordinated}

\section{Certification à profondeur \texorpdfstring{\(10^4\)}{10^4} des développements coinductifs de \texorpdfstring{\(\pi\), \(e\) et \(\varphi\)}{pi, e et phi}}
\label{sec:truncamiento-10000-pi-e-phi-rev9}

\noindent\textbf{Portée mathématique.}
Le théorème uniforme précédent quantifie sur tout entier positif \(N\) et
construit l'objet infini comme limite inverse de ses cylindres compatibles.
Ce chapitre exécute le cas \(N=10^4\) comme certificat fini particulièrement
exigeant ; \(10^4\) n'est pas une borne du générateur, de même qu'une exécution
à un million de chiffres serait une autre instance du même algorithme et non une
modification de sa loi.

La présente exécution fixe une profondeur de calcul et vérifie que le
cylindre obtenu est contenu dans une unique cellule décimale de dix mille
positions. Le plus petit nombre de blocs compatible avec la monodromie nonadique
et avec
\[
  3^{6B}>10^{10000}
\]
est
\[
  B=3501=389\cdot9,
  \qquad 6B=21006\ \text{trits}.
\]

Pour chacun des trois canaux, les \(729\)
sous-cylindres sont parcourus exhaustivement à chaque étape et l'on conserve l'unique sous-cylindre compatible avec l'encadrement
rationnel du lecteur. Le calcul produit
\[
 \begin{array}{c|r}
 \text{grandeur calculée par canal}&\text{valeur}\\
 \hline
 \text{partitions complètes }729\to1&3501\\
 \text{trits transportés}&21006\\
 \text{paires de même phase }(K,K+9)&3492\\
 \text{tours nonadiques complets}&388
 \end{array}
\]
et ne nécessite pas de nonade supplémentaire. Les \(396\) premiers blocs
coïncident exactement avec le témoin précédent à mille chiffres.

\Needspace{28\baselineskip}
Les lecteurs de clôture, de propagation et d'auto-échelle sont évalués au moyen
d'inégalités rationnelles exactes. L'accélération du calcul réorganise les
sommes par l'arithmétique entière et n'ouvre ni bibliothèques de constantes ni préfixes
décimaux cibles. Ce n'est qu'après avoir écrit les sorties que l'on exécute un
second calcul indépendant. On vérifie, dans cet ordre, la
génération rationnelle des trois canaux, la composition G9--T1 de
\(\alpha_{\rm HMT}\) et la coïncidence des dix mille chiffres avec les
développements de référence consultés seulement à la fin. Les trois comparaisons
sont exactes sur les dix mille positions calculées. Cette vérification finie
ne remplace pas l'argument uniforme de prolongation à profondeur arbitraire.

La démonstration étendue précise aussi le sens du retour. Dans les
\(3492\) paires par canal, la phase est conservée et l'état
complet change toujours, car la profondeur, le préfixe et le cylindre avancent. Une projection locale
réduite peut se répéter dans certaines paires sans que l'état TPK revienne. Cette
distinction remplace l'affirmation plus forte d'apériodicité locale par la
propriété effectivement démontrée : monodromie de phase avec mémoire de l'état
enrichi.

Les quatre développements complets sont imprimés dans
\cref{app:desarrollos-10000-rev9}. Les récurrences, les intervalles
rationnels et les comptages précédents déterminent le calcul ; les développements
imprimés sont leurs troncatures finies et permettent de comparer chaque chiffre sans
devenir des prémisses du générateur.

\chapter{Certificats de précision arbitraire et échantillons décimaux finis reproductibles}
\label{app:desarrollos-10000-rev9}

Cet appendice imprime les sorties complètes de l'exécution décrite dans
\cref{sec:truncamiento-10000-pi-e-phi-rev9,sec:alpha-t1-10000-rev9}. Les
sauts de ligne sont exclusivement typographiques. Chaque développement constitue
un témoin fini de la loi coinductive de précision arbitraire démontrée dans
la Partie III ; il ne définit pas la constante ni ne limite la profondeur du générateur.

\begingroup
\lstset{basicstyle=\ttfamily\tiny,breaklines=true,breakatwhitespace=false,
  frame=single,columns=fullflexible,keepspaces=true,numbers=none}

\section{Développement coinductif de \(\pi\) et certificat de troncature compatible}
\lstinputlisting{generated/decimals_10000_rev9/pi_10000_decimales_96col.txt}

\section{Développement coinductif de \(e\) et certificat de troncature compatible}
\lstinputlisting{generated/decimals_10000_rev9/e_10000_decimales_96col.txt}

\section{Développement coinductif de \(\varphi\) et certificat de troncature compatible}
\lstinputlisting{generated/decimals_10000_rev9/phi_10000_decimales_96col.txt}

\section{Développement de \texorpdfstring{\(\alpha_{\rm HMT}\)}{alpha HMT}}
\lstinputlisting{generated/decimals_10000_rev9/alpha_T1_10000_decimales_96col.txt}
\endgroup


\chapter{Formalisation assistée du noyau algébrique}
\label{app:formalizacion-asistida-acumulativa}

La formalisation vérifie les identités finies sur les matrices, les algèbres et les
états dodécaphasiques typés. Sa fonction est d'isoler les obligations
algébriques susceptibles de réduction mécanique et de maintenir séparées la
provenance des données, l'énumération finie et la preuve formelle.

\HMTDemoteHeadings
\chapter{Vérification formelle d'identités algébriques finies}
\label{app:lean}

Les égalités finies suivantes sont formalisées dans Lean et décidées par
réduction du noyau logique, sans axiomes externes supplémentaires. Les théorèmes
démontrent :
\begin{longtable}{p{.40\textwidth}p{.50\textwidth}}
\toprule
Théorème & Contenu\\
\midrule
\endhead
\texttt{aw\_symmetric}
& symétrie de la matrice de Paley--Witt \(A_W\) ;\\
\texttt{aw\_square\_minus\_identity}
& \(A_W^2=-I_6\) sur \(\mathbb F_3\) ;\\
\texttt{x\_sq\_add\_one\_has\_no\_root}
& irréductibilité de \(x^2+1\) sur \(\mathbb F_3\) ;\\
\texttt{dual\_has\_nonzero\_nilpotent}
& existence d'un nilpotent dans l'algèbre parabolique ;\\
\texttt{split\_has\_nontrivial\_idempotent}
& existence d'un idempotent non trivial dans l'algèbre hyperbolique ;\\
\texttt{canonical\_charge}
& \(\sum_{m=1}^{12}K_m=6263\);\\
\texttt{hadamard\_orbit\_one/two/three}
& les trois orbites de Hadamard produisent l'état dodécaphasique publié.\\
\bottomrule
\end{longtable}

La formalisation reçoit les matrices et le vecteur dodécaphasique comme données
typées et vérifie leurs conséquences algébriques. La construction de l'état
signé et l'énumération de (104\,976) germes appartiennent aux théorèmes
combinatoires antérieurs et ne sont pas remplacées par cette réduction formelle.

\HMTRestoreHeadings

\chapter{Correspondance entre énoncés, preuves, certificats et propriétaires scientifiques}
\label{app:indice-resultados-fuentes-acumulativo}

Le tableau suivant identifie la résidence mathématique principale de chaque
résultat et le contrôle négatif qui distingue sa dérivation d'une simple
coïncidence. Aucune ligne ne remplace l'énoncé ni la démonstration principale.

\begingroup
\small
\setlength{\tabcolsep}{4pt}
\renewcommand{\arraystretch}{1.16}
\begin{longtable}{@{}>{\raggedright\arraybackslash}p{.20\textwidth}
  >{\raggedright\arraybackslash}p{.43\textwidth}
  >{\raggedright\arraybackslash}p{.27\textwidth}@{}}
\caption{Résultats, démonstrations de référence et contrôles reproductibles.}
\label{tab:indice-fuentes-acumulativo}\\
\toprule
Résultat & Démonstration de référence & Control reproducible\\
\midrule
\endfirsthead
\toprule
Résultat & Démonstration de référence & Control reproducible\\
\midrule
\endhead
Construction de l'APP, double évaluation et saut spectral
& Chapitres de l'APP et de l'arithmétique généalogique
& Identités de support, spectre et reconstruction des deux feuillets\\
État TRIT, orientation et trois algèbres quadratiques
& Chapitre du TRIT et classification algébrique
& Conservation du résidu, du quotient, de l'orientation et de la retenue\\
Inventaire opératoire complet du TPK
& Chapitres de transport, mémoire et opérateurs fondamentaux
& Commutativité, naturalité et prolongement compatible\\
Extrême et moyenne raison holographique
& Complétion cubique nonadique
& Identité $1000=729+243+27+1$ et conservation de l'unité\\
Génération prolongeable de \(\pi,e,\varphi\)
& Biographies coinductives et monodromie nonadique
& Cylindres emboîtés, troncatures compatibles et falsificateur d'oracle\\
Deux constructions internes de \(\alpha_{\rm HMT}\)
& Roue dodécaphasique et cartes réversibles
& Équivalence des sorties et séparation de la reconnaissance CODATA\\
Projecteurs dodécaphasiques et défaut \(-1/12\)
& Noyau Hensel--Witt
& Symétrie, idempotence, rangs $3/4$ et trace normalisée\\
Double projection archimédienne et exceptionnelle
& Lecteur de cylindres et lecteur d'incidence
& Même état d'entrée, codomaines typés et non-identification des sorties\\
Réseau de Leech, VOA et Moonshine
& Corridor exceptionnel complet
& Opérateurs d'entrelacement d'incidence et obstructions d'équivariance\\
Généalogie du nombre et fonctionnelles spectrales
& Clôture du concept de nombre et arithmétique spectrale
& Produits d'Euler, caractères et bornes dirigées\\
Droite d'action et ordre décimal \(-34\)
& Action bidirectionnelle et coordonnées angulaires conjuguées
& Conoyau d'ordre $16$ et séparation du correcteur relatif\\
Barbero--Immirzi, opérateur d'aire et spectre
& Passage hexade--octade et hamiltonien modulaire
& Canaux $90/120$, aire, première loi et entropie relative\\
Angles CKM et violation CP
& Chapitre du mélange et de la réflexion rationnelle
& Unitarité, Jarlskog et prudence concernant la réflexion physique\\
Électron bidirectionnel et masse raffinée
& Définition de particule et première réalisation électronique
& Point fixe central, lecteurs $K_D/K_\Omega$ et valeur bêta directrice\\
Atlas \(81/56/13/324\) et opérateurs de masse
& Loi générale des masses et atlas route par route
& Couverture complète, secteur nul et sélection prospective\\
Permittivité et perméabilité du vide
& Réponse constitutive bicouche
& Reconstruction séparée de \(Z,\mu,\varepsilon,c\)\\
Formulation variationnelle et hamiltonienne de la gravité
& Gravitation Einstein--Cartan--Holst
& Équivalence Euler--Lagrange/Hamilton et contrôles torsionnels\\
Navier--Stokes et Yang--Mills
& Clôtures terminales des réalisations fluide et de jauge
& Matrice PF84, coercivité et transfert\\
\bottomrule
\end{longtable}
\endgroup

Les formulations réduites précédentes ne gouvernent pas les théorèmes de la
Partie IV. La portée de chaque spécialisation y est déterminée par son objet,
son domaine, ses applications et ses identités de comparaison. En particulier, une
étiquette historique ne peut à elle seule déclarer une conclusion ouverte ou
fermée, de même qu'une classification documentaire ne remplace pas la
commutativité d'un diagramme. Les constructions finies du tableau précédent
sont conservées comme antécédents ; leur prolongement est décidé au moyen des
comparateurs explicites de chaque chapitre.

\chapter{Domaines, cartes et interfaces de réalisation}
\label{app:dominios-cartas-interfaces-acumulativo}

Ce bloc sépare les théorèmes internes, les cartes d'évaluation, les
interfaces physiques et les programmes d'extension. Cette séparation permet de lire
chaque égalité dans son domaine précis et maintient visibles les données, les jauges,
les orientations ou les conditions de frontière qui interviennent dans une réalisation.

\HMTDemoteHeadings
\chapter{Domaines, codomaines et hypothèses des résultats}
\label{app:premisas}

\begin{enumerate}
\item Le recensement \(104\,976\to468\to243\), l'image de 42 triplets, la
factorisation \(468=26\times18\), les cinq régions et les cinq cycles sont des
théorèmes exhaustifs du système fini de transport spécifié dans cet
ouvrage.

\item Le catalogue enrichi produit, par des prédicats finis de fibre et sans
consultation décimale, l'unique orbite de clôture, l'unique orbite de
propagation et l'unique orbite d'auto-échelle ; le calibre interne y oriente
les représentants \(010211,201101,121200\). Leurs caractères sont
construits, respectivement, par la clôture pentarégionale avec retour
unitaire, la récurrence de semi-groupe \((n+1)a_{n+1}=a_n\) et l'unique rayon positif
de l'incidence \(F_{\rm av}\). L'évaluation archimédienne publie
alors les valeurs engendrées et la reconnaissance conventionnelle les nomme
\(\pi,e,\varphi\). Pour chaque profondeur \(N\), le prolongement TPK produit
un préfixe unique compatible avec celui de profondeur \(N-1\) ; la limite
inverse de ces préfixes fixe le développement complet. Ainsi, les mots,
les caractères et les valeurs sont des étapes successives d'une construction
interne, non des étiquettes choisies au moyen de tables externes.

\item Le théorème dodécaphasique a pour domaine causal l'état terminal
enrichi \(x_{\rm term}\) produit par APP--TRIT--TPK. La carte signée
publie à partir de lui
\(x_{\rm term}\xrightarrow{R_{\rm sgn}}U_{12}^{\rm sgn}\), et l'inversion
entière publie ensuite
\(U_{12}^{\rm sgn}\xrightarrow{\mathcal H_{12}^{-1}}K\). Dans ce domaine,
les correspondances avec \((D_3K,D_4K,Q)\), la représentation de Hadamard et
la clôture dodécaphasique sont exactes et réversibles. Ni \(U\) ni \(K\) ne sont des
entrées génératrices. La trace réduite d'une trajectoire positionnelle et le
système bivarié complet sont des objets distincts, comme l'établit le
\cref{chap:tpk-definicion}.

\item L'égalité avec la chaîne centrale 2018 est une reconnaissance externe.
Le mot mathématique interne et sa démonstration sont indépendants des
recommandations métrologiques ultérieures.

\item L'instance finie de la double projection démontre un drapeau commun,
\(W_{12}\) et \(M_{12}\). Dans la branche réticulaire, on fixe le recollement de
\(A_2^{12}\), la chambre radiale et l'origine ; le 3-voisin résultant est certifié
pair, unimodulaire et sans racines, et est identifié classiquement à Leech. La
construction de \(V^\natural\) et du Monstre est une conséquence classique. La
flèche \(W_{12}\to W_{24}\) est réalisée au moyen du plan résiduel d'une
dodécade et du relèvement unique de ses hexades, relativement à un choix
\((D,\ell)\) ; les identifications s'effectuent au moyen d'applications, non
par des égalités de cardinalité. La compatibilité projective de ces
données et leur double réalisation sur l'état enrichi commun sont démontrées
dans les
\cref{p12-83-c17-s2:thm:lector-excepcional-global-rev6,p12-83-c17-s2:thm:doble-realizacion-global-rev6}.

\item Le théorème des cylindres produit un unique réel pour chaque histoire de son
domaine. La loi de prolongement couvre de manière cohérente chaque compact et
l'exhaustion dirigée couvre la droite entière, comme le démontrent les
\cref{p12-83-c18-s1:thm:sobreyectividad-evaluacion-compacto,p12-83-c18-s1:prop:agotamiento-recta-real-u024,p12-83-c18-s2:hmtch:thm-cociente-real}.
Par conséquent, l'évaluation est surjective et son quotient archimédien est
\(\mathbb R\) ; la couverture est une conclusion du générateur HMT, non une
prémisse supplémentaire.

\item Le produit natif est démontré dans le secteur canonique des mots.
L'extension monoïdale à l'espace complet des histoires et l'entrelacement
avec les filtrations de survie sont formulés comme des applications distinctes
dans le \cref{chap:producto-nativo}.
\end{enumerate}

\chapter{Invariance causale des interfaces par rapport aux références externes}
\label{chap:independencia-prospectiva}
\index{indépendance prospective}\index{graphe causal}
\index{protocolo ciego}\index{tours!complétude reconstructive}

\section{Critère d'indépendance causale par rapport aux références externes : graphe dirigé acyclique et fixation préalable de la sortie}
\label{sec:definiciones-independencia-prospectiva}

Soient \(\mathsf P\) les primitives internes, \(\mathsf R\) les règles,
lecteurs, calibres et normalisations déclarés, \(G\) un générateur,
\(o=G(\mathsf P,\mathsf R)\) sa sortie et \(y\) une référence externe.

\begin{definition}[Interface prospectivement indépendante]
\label{def:interfaz-prospectivamente-independiente}
L'interface est prospectivement indépendante de \(y\) si :
\begin{enumerate}
 \item \(y\notin\operatorname{Anc}(o)\);
 \item \(\mathsf P,\mathsf R,G\) et les contrôles négatifs sont figés avant
 l'ouverture de \(y\) ;
 \item l'exécution n'a accès ni directement ni indirectement à \(y\) ;
 \item \(o\) et sa trace sont scellés avant le calcul d'une comparaison
 \(\operatorname{cmp}(o,y)\).
\end{enumerate}
\end{definition}

\begin{definition}[Indépendance de la sélection]
\label{def:independencia-seleccion}
Le caractère aveugle de chaque exécution ne suffit pas si, après avoir produit des candidats
\(G_j(\mathsf P,\mathsf R_j)\), on publie
\[
 j_\ast(y)=\arg\min_j
 \operatorname{cmp}\!\left(G_j(\mathsf P,\mathsf R_j),y\right).
\]
L'indépendance de la sélection interdit que cette opération détermine la sortie
publiée.
\end{definition}

\begin{definition}[Prédiction multivaluée]
\label{def:prediccion-multivaluada}
Si les descripteurs physiques non métrologiques et les restrictions internes
admettent plusieurs multisections, la sortie prospective est l'ensemble complet
\[
 \widehat{\mathfrak R}(d)
 =\{\mathcal M:\mathcal M\text{ satisfait toutes les contraintes
 préinscrites pour }d\}.
\]
La non-unicité est une sortie de la procédure. Choisir ensuite un élément
par proximité d'une grandeur est interdit.
\end{definition}

\section{Théorème de non-interférence}
\label{sec:teorema-no-interferencia}

\begin{definition}[Graphe causal dirigé acyclique fidèle et complet]
\label{def:dag-fiel-completo}
Un graphe dirigé acyclique \(D\) est fidèle si chaque nœud représente la même
fonction et les mêmes entrées que la construction citée. Il est complet pour
une sortie si toute dépendance causalement pertinente apparaît comme arête ou
racine déclarée. La fidélité et la complétude sont des hypothèses sémantiques ;
l'acyclicité de l'ensemble de données ne les démontre pas.
\end{definition}

\begin{theorem}[Indépendance prospective conditionnée par le graphe causal]
\label{thm:no-interferencia-prospectiva}
Soit un système déterministe d'équations structurelles représenté par un
graphe dirigé acyclique fini \(D\), fidèle et complet pour \(o\). Si
\(y\notin\operatorname{Anc}_D(o)\), toutes les racines ancestrales de \(o\)
restent figées et l'exécution n'a pas d'accès caché à \(y\), alors
pour tous \(y_0,y_1\),
\[
 o_{\operatorname{do}(y=y_0)}
 =o_{\operatorname{do}(y=y_1)}.
\]
\end{theorem}

\begin{proof}
On ordonne topologiquement les antécédents de \(o\). Leurs racines sont internes
et restent fixes. Par induction sur l'ordre, chaque nœud reçoit les mêmes
arguments sous les deux interventions et, par déterminisme, renvoie la même
valeur. L'égalité se propage jusqu'à \(o\). L'empreinte cryptographique ne crée pas
l'indépendance ; elle rend seulement détectable une modification ultérieure.
\end{proof}

\begin{corollary}[Permutation de la comparaison externe]
\label{cor:permutacion-comparacion-externa}
Si les noms, les lignes ou les valeurs externes sont permutés après le scellement, la
sortie interne ne change pas, bien que son score puisse changer :
\[
 o_{\operatorname{do}(y=\sigma y_0)}
 =o_{\operatorname{do}(y=y_0)},\qquad
 \operatorname{cmp}(o,\sigma y_0)
 \ne\operatorname{cmp}(o,y_0).
\]
\end{corollary}

\begin{proposition}[Séparation des branches calibrées]
\label{prop:separacion-ramas-calibradas}
Si une référence externe est intervenue dans une sélection historique, la
dépendance ne contamine pas les objets internes qui ne descendent pas de cette
sélection. La branche peut être isolée et soumise à un nouveau protocole prospectif
sans modifier le noyau.
\end{proposition}

\begin{proof}
Être un antécédent est une propriété locale du graphe. Les branches
\[
 Z_0\longrightarrow p=5\longrightarrow\text{échelle électronique},
\qquad
 \{\text{masses}\}\longrightarrow\{\text{signatures calibrées}\}
 \longrightarrow\text{table}
\]
restent hors du sous-graphe protégé. Les tours, le caractère formel,
l'atlas et le centre \((5,5)\) ne reçoivent pas d'arêtes de ces références.
\end{proof}

\section{Invariance contrefactuelle sur le graphe causal dirigé et portée de la conclusion}
\label{sec:certificado-independencia-prospectiva}

Le graphe causal publié contient \(71\) nœuds, \(89\) arêtes,
\(16\) interfaces et
\(29\) sorties protégées. La vérification causale établit l'acyclicité, les décomptes,
les voies déclarées, l'existence des sources enregistrées, l'absence de
références externes parmi les antécédents consignés et l'invariance
contrefactuelle du modèle symbolique. Le résultat est
\[
 \#\{\text{prédécesseurs externes de sorties protégées}\}=0.
\]
Le certificat ne démontre pas à lui seul qu'aucune arête ne manque et n'audite pas
toutes les lectures de données de chaque programme scientifique. Sa force est donc
\textsc{démontré avec une structure de départ explicite} : la conclusion
s'obtient sur le graphe dirigé acyclique expressément déclaré et reste
conditionnée par sa fidélité et sa complétude.

Le remplacement simultané des références CODATA, Planck,
Barbero--Immirzi, CKM, masses, gravité et cristal temporel peut modifier les
comparaisons et les branches calibrées. Il ne modifie pas, dans le graphe déclaré,
le catalogue de \(468\) émissions, les cinq régions de \(\pi\), la
génération coinductive corrélée, \(\alpha_{12}\), \(C^\ast\), la décade \(-34\), les tours,
l'atlas, le centre électronique, la double projection ni le corridor
exceptionnel.

\section{Essai multivalué sans accès à la valeur cible et critère d'identification}
\label{sec:protocolo-ciego-multivaluado}

\begin{algorithm}[Attribution future sans grandeurs]
\label{alg:asignacion-futura-sin-magnitudes}
L'interface physique s'exécute en quatre phases irréversibles :
\begin{enumerate}
 \item \emph{Gel} : on scelle l'autorité, l'atlas, les involutions,
 les restrictions, le code, la priorité et les contrôles négatifs.
 \item \emph{Entrée autorisée} : seuls sont reçus le secteur, la génération, la charge,
 la couleur, la conjugaison, le spin/la statistique, la composition et le type métrologique ;
 les masses, les signatures raffinées, les résidus, les angles et les unités sont exclus.
 \item \emph{Sortie} : toutes les multisections compatibles sont énumérées ; pour
 chaque extension de chaque multisection, on calcule la voie observable, le registre,
 la signature, \(\chi_{\rm form}\), la spécialisation et l'observable. Dans les objets
 composés, on raccorde d'abord les extensions et les registres. On scelle le
 multiensemble complet.
 \item \emph{Ouverture} : un dépositaire distinct ouvre les grandeurs et calcule
 les résidus, les incertitudes et les covariances sans modifier aucun candidat.
\end{enumerate}
\end{algorithm}

La multisection enrichie, et non une extension ponctuelle sélectionnée
rétrospectivement, est la sortie de l'appariement physique. L'évaluation
ultérieure peut être effectuée fibre par fibre. L'algorithme précédent est une
spécification réfutable ; il n'est pas présenté comme une prédiction multiensemble déjà
exécutée.

\begin{criterion}[Échec du protocole en aveugle]
\label{crit:fallo-protocolo-ciego}
Le protocole échoue si le code peut lire une grandeur, si une signature
historique calibrée réapparaît comme entrée, si un candidat est écarté après
ouverture de la référence, si une masse positive est confondue avec une masse exactement nulle
ou si le multiensemble scellé change lorsque le catalogue externe est permuté.
\end{criterion}

\section{Caractère formel et complétude reconstructive}
\label{sec:completitud-reconstructiva-torres}

Pour \(v=(n_A,s_C,k,b_{120},b_{\zeta})\in\ZZ^5\), le caractère formel
antérieur à toute spécialisation est
\[
 \chi_{\rm form}(v;x_A,x_C,x_\Delta,x_{120},x_{270})
 =\exp\!\left(
 n_Ax_A+s_Cx_C+kx_\Delta
 +b_{120}x_{120}+b_{\zeta}x_{270}\right).
\]
La spécialisation déclarée
\[
 \Theta_0
 =\left(A_{\rm rad},\frac{C^\ast_{\rm rad}}6,
 \Dfour,s_{120},\zeta_{270}\right),
 \qquad \zeta_{270}=135\Delta_4^2,
\]
produit
\(\chi_0(v)
=\chi_{\rm form}(v;\Theta_0)\).
La signature brute, le caractère formel, sa spécialisation centrale, le
raffinement global \(\eta\in\mathcal E_{\mathfrak S}\) et un tableau
calibré sont cinq objets différents.

\begin{theorem}[Complétude reconstructive positive]
\label{thm:completitud-reconstructiva-positiva}
Soit
\[
 \mathcal N=\langle90,120\rangle_{>0}
\]
et soit \((n_j)_{j\ge1}\) son énumération croissante. Étant donnés
\(0<q_+<q_-<1\), définissons
\[
 a_j=q_-^{n_j}-q_+^{n_j}>0.
\]
Alors \(a_j\to0\) et \(\sum_ja_j<\infty\). Une règle déterministe étant fixée
pour départager les demi-entiers, pour tout \(r\in\RR\), la récurrence
\[
 r_0=r,\qquad
 \nu_j=\operatorname{nint}\!\left(\frac{r_{j-1}}{a_j}\right),\qquad
 r_j=r_{j-1}-\nu_ja_j
\]
satisfait
\[
 |r_j|\leq\frac{a_j}{2}\longrightarrow0,\qquad
 r=\sum_{j\ge1}\nu_ja_j,\qquad
 \sum_{j\ge1}|\nu_ja_j|<\infty.
\]
Ainsi, tout \(R>0\), avec \(r=\log R\), est représenté exactement dans la
complétion par
\(\exp(\sum_j\nu_ja_j)\).
\end{theorem}

\begin{proof}
Les indices de \(\mathcal N\) sont \(90,120\) et tous les multiples de \(30\)
à partir de \(180\). Par conséquent,
\[
 \sum_{n\in\mathcal N}q_-^n
 =q_-^{90}+q_-^{120}
 +\frac{q_-^{180}}{1-q_-^{30}}<\infty.
\]
Comme \(0<a_j<q_-^{n_j}\), la suite est sommable et converge vers zéro. La
définition de l'entier le plus proche implique
\(|r_j|\le a_j/2\). Par télescopage,
\(\sum_{i=1}^{N}\nu_ia_i=r-r_N\to r\). Pour \(j\ge2\),
\[
 |\nu_ja_j|=|r_{j-1}-r_j|
 \le\frac{a_{j-1}+a_j}{2},
\]
et la sommabilité de \((a_j)\) prouve la convergence absolue.
\end{proof}

Le théorème est une complétude du dictionnaire sur \(\RR_{>0}\). Si les
coefficients \(\nu_j\) sont obtenus à partir du résidu d'une masse connue, leur utilisation
est une \textsc{reconstruction calibrée} ; elle ne fournit pas la sélection
prospective. Elle ne représente pas non plus une masse exactement nulle. Le photon, le gluon ou un
éventuel mode de spin deux exigent une protection exacte par symétrie dans la
représentation physique.

\Needspace{34\baselineskip}
\section{Protocoles prospectifs d'invariance pour les familles de constantes}
\label{sec:protocolos-prospectivos-especificos}

\begin{description}
 \item[Barbero--Immirzi.] On fige \(S_{90},S_{120}\), les poids
 \((12,-1)\) et leur normalisation avant d'ouvrir une référence physique.
 L'évaluation de \(\Gamma_{6\to8}\) est aveugle à cette référence,
 mais l'unicité préalable des séries et des poids et l'opérateur d'entrelacement avec l'opérateur
 d'aire sont des problèmes distincts.
 \item[CKM et CP.] On scelle les quatre équations, leurs coefficients et la
 convention angulaire ; on publie simultanément les quatre paramètres,
 \(|V|\) et \(J\). La confrontation utilise une covariance fixée à l'avance.
 \item[Constitutives.] Les identités adimensionnelles précèdent
 \(Z_{\rm ref}\) et \(c_{\rm ref}\). Le changement de carte transforme la
 réalisation, non les identités, et ne crée pas quatre prédictions
 indépendantes.
 \item[Gravité et inertie.] Les rapports
 \(G_{\rm tor}/I_{\rm tor}=1\) et
 \(G_{\rm av}/I_{\rm av}=\pi/\varphi^2\) appartiennent à des couches différentes.
 Leur interprétation physique exige un morphisme et une confrontation prospective ; un
 phénomène cosmologique ne sélectionne pas le rapport.
\end{description}

\section{Stratification probatoire : identités internes, déductions, reconstructions calibrées et interfaces soumises à une fermeture mathématique}
\label{sec:cinco-niveles-independencia}

\begin{table}[htbp]
\centering
\caption{Niveaux que l'indépendance causale ne permet pas de confondre.}
\label{tab:cinco-niveles-independencia}
\small
\begin{tabularx}{\textwidth}{@{}L{0.25\textwidth}YL{0.27\textwidth}@{}}
\toprule
Niveau & Exemple & Énoncé autorisé\\
\midrule
Générateur interne certifié&
Régions, \(\alpha_{12}\), \(C^\ast\), \(-34\), atlas&
Sortie obtenue dans le domaine déclaré.\\
Évaluation conditionnée en aveugle&
\(\Gamma_{6\to8}\), CKM fixé, \(\chi_0(v)\)&
Ne consulte pas la cible, mais dépend de la fonctionnelle ou de la signature.\\
Protocole en aveugle non exécuté&
Spécification multivaluée&
Il existe une procédure réfutable ; non une sortie scellée.\\
Reconstruction calibrée&
Signatures historiques et raffinement à partir de \(r=\log R\)&
Évaluation exacte conditionnée par une référence.\\
Interface physique en attente&
Particule--antiparticule, aire, masse nulle, gravité&
Domaine défini ; représentation ou confrontation manquante.\\
\bottomrule
\end{tabularx}
\end{table}

\HMTRestoreHeadings

\chapter{Reproductibilité de l'atlas et comparaison des réalisations de masse}
\label{app:reproducibilidad-masas-acumulativa}
\label{app:suplemento-digital-masas-rev11}

La théorie des masses est engendrée sur le domaine complet de
\(81\) positions, \(56\) familles de routes, \(13\) multisections,
\(324\) routes orientées et le secteur nul. La phase structurelle calcule
\[
\begin{aligned}
&\text{position ou fibre}\longrightarrow\text{famille}
\longrightarrow\text{route orientée}\\
&\longrightarrow\text{registre}
\longrightarrow\text{signature nominale}\longrightarrow\text{caractère et tour},
\end{aligned}
\]
et, lorsque la fibre possède le lecteur opératoriel correspondant, ajoute la
branche de raffinement
\[
 \text{registre}\longrightarrow(K_D,K_\Omega)
 \longrightarrow\text{spectre de fibre}.
\]
La seconde branche raffine la première et ne la remplace pas : en particulier,
\(W\), \(Z\) et \(H\) conservent leurs routes CT--108, leurs signatures angulaires et
leurs caractères scalaires bien qu'ils n'appartiennent pas à l'atlas opératoriel de
treize multisections. Ces invariants sont fixés avant la confrontation externe ;
une référence métrologique n'intervient pas dans la sélection de la route, de la position,
de la signature ni du coefficient.

\section{Invariants finis de la classification}
\label{sec:fuentes-certificado-tabla-masas-acumulativa}

\begin{itemize}
 \item atlas complet de $81$ cellules avec leurs lecteurs ;
 \item $56$ familles de routes avec leurs spectres ;
 \item $13$ opérateurs de multisection ;
 \item $324$ routes orientées avec registre et caractère ;
 \item inventaires complets des masses et des largeurs ;
 \item carte comparative postérieure à la sélection généalogique.
\end{itemize}

% REMATE_LOCAL_AP_I
\Needspace{25\baselineskip}
\section{Conditions d'exactitude, de stabilité et de reproductibilité des certificats}
\label{sec:ejecucion-tabla-masas-acumulativa}

Le contrôle intégral exige la couverture \(81/56/13/324\), le secteur nul, les
\(14\) profils de \(K_D\), les \(9\) profils de \(K_\Omega\), les
\(40\) paires conjointes et les \(18\) correspondances de lecteurs. Il vérifie
en outre que la masse électronique raffinée est évaluée à partir d'une unique expression,
que les secteurs nuls conservent \(\kappa\) non applicable, qu'une ligne
opératoire ne publie pas de résiduel scalaire, que les douze routes nominales
\(e,\mu,\tau,u,d,s,c,b,t,W,Z,H\) conservent leur caractère scalaire et que la
réalisation physique maintient le codomaine propre de chaque sortie. L'absence d'une paire
\(K_D/K_\Omega\) dans un secteur ne peut dégrader ni effacer son lecteur scalaire
natif.

La sortie électronique reproduite est
\[
m_e^\beta
=m_e^{(0)}R_{\rm act}^{80-54\alpha_{\rm HMT}+6\Delta_4}
=0.510998950690000808608257165\ldots\;\mathrm{MeV}/c^2.
\]
Les fibres de rang supérieur publient leur paire opératoire, leurs spectres et le
morphisme de réalisation correspondant. La sélection des valeurs propres est
déterminée par ce morphisme avant la comparaison externe.


\section{Domaine scientifique de la classification}
\label{sec:alcance-publico-apendices-acumulativo}

Le domaine des masses est l'atlas complet précédent. La construction reçoit
exclusivement des positions, des familles, des multisections, des routes, des registres et des
opérateurs typés ; les données de provenance demeurent séparées du contenu
scientifique.
